\documentclass[twoside,11pt]{article}

%+
%  Name:
%     sun232.tex

%  Purpose:
%     SUN documentation for ORAC-DR imaging (SUN/232)

%  Notes:
%     -  Has definitions to give small capitals for CCDPACK,
%     CONVERT, CURSA, GAIA, KAPPA, PHOTOM, and PISA.  Each
%     command is \package.
%     -  Definition of \ORACDR as ORAC-DR in small capitals.
%     -  Definition of \FITSref to point to FITS Home Page.
%     -  SST has singled-sided name.
%     -  Added \pagebreak[3] and \goodbreaks to SST to prevent
%     section headings at the foot of a page.
%     -  Used SUN/95 version of sstitemlist to remove
%     superfluous space after headings.  Added ssthitemlist
%     where the spacing is correct.
%     -  Uses some definitions of degrees, arcminutes,
%     arcsecond, and underscore.

%  Authors:
%     MJC: Malcolm J. Currie (JACH)

%  History:
%     2000 February 2 (MJC):
%     �  Original version 232.0.

%  Copyright:
%     Copyright (C) 1998--2000 Particle Physics and Astronomy
%     Research Council.  All Rights Reserved.

%-

% ? Specify used packages
\usepackage{graphicx}        %  Use this one for final production.
%\usepackage[draft]{graphicx} %  Use this one for drafting.
% ? End of specify used packages

\pagestyle{myheadings}

% -----------------------------------------------------------------------------
% ? Document identification
% Fixed part
\newcommand{\stardoccategory}  {Starlink User Note}
\newcommand{\stardocinitials}  {SUN}
\newcommand{\stardocsource}    {sun\stardocnumber}

% Variable part - replace [xxx] as appropriate.
\newcommand{\stardocnumber}    {232.0}
\newcommand{\stardocauthors}   {Malcolm J. Currie\\
                               Joint Astronomy Centre, Hilo, Hawaii}
\newcommand{\stardocdate}      {2000 February 2}
\newcommand{\stardoctitle}     {ORAC-DR - imaging data reduction}
\newcommand{\stardocversion}   {1.0-0}
\newcommand{\stardocmanual}    {User Guide}
\newcommand{\stardocabstract}  {{\footnotesize ORAC-DR} is a
general-purpose automatic data-reduction pipeline environment.  This
document describes its use to reduce imaging data collected at the
United Kingdom Infrared Telescope (UKIRT) with the UFTI and IRCAM
instruments.  It outlines the algorithms used and how to make minor
modifications of them, and how to correct for errors made at the
telescope.}

% ? End of document identification
% -----------------------------------------------------------------------------

% +
%  Name:
%     sun.tex
%
%  Purpose:
%     Template for Starlink User Note (SUN) documents.
%     Refer to SUN/199
%
%  Authors:
%     AJC: A.J.Chipperfield (Starlink, RAL)
%     BLY: M.J.Bly (Starlink, RAL)
%     PWD: Peter W. Draper (Starlink, Durham University)
%
%  History:
%     17-JAN-1996 (AJC):
%        Original with hypertext macros, based on MDL plain originals.
%     16-JUN-1997 (BLY):
%        Adapted for LaTeX2e.
%        Added picture commands.
%     13-AUG-1998 (PWD):
%        Converted for use with LaTeX2HTML version 98.2 and
%        Star2HTML version 1.3.
%     {Add further history here}
%
% -

\newcommand{\stardocname}{\stardocinitials /\stardocnumber}
\markboth{\stardocname}{\stardocname}
\setlength{\textwidth}{160mm}
\setlength{\textheight}{230mm}
\setlength{\topmargin}{-2mm}
\setlength{\oddsidemargin}{0mm}
\setlength{\evensidemargin}{0mm}
\setlength{\parindent}{0mm}
\setlength{\parskip}{\medskipamount}
\setlength{\unitlength}{1mm}

% -----------------------------------------------------------------------------
%  Hypertext definitions.
%  ======================
%  These are used by the LaTeX2HTML translator in conjunction with star2html.

%  Comment.sty: version 2.0, 19 June 1992
%  Selectively in/exclude pieces of text.
%
%  Author
%    Victor Eijkhout                                      <eijkhout@cs.utk.edu>
%    Department of Computer Science
%    University Tennessee at Knoxville
%    104 Ayres Hall
%    Knoxville, TN 37996
%    USA

%  Do not remove the %begin{latexonly} and %end{latexonly} lines (used by 
%  LaTeX2HTML to signify text it shouldn't process).
%begin{latexonly}
\makeatletter
\def\makeinnocent#1{\catcode`#1=12 }
\def\csarg#1#2{\expandafter#1\csname#2\endcsname}

\def\ThrowAwayComment#1{\begingroup
    \def\CurrentComment{#1}%
    \let\do\makeinnocent \dospecials
    \makeinnocent\^^L% and whatever other special cases
    \endlinechar`\^^M \catcode`\^^M=12 \xComment}
{\catcode`\^^M=12 \endlinechar=-1 %
 \gdef\xComment#1^^M{\def\test{#1}
      \csarg\ifx{PlainEnd\CurrentComment Test}\test
          \let\html@next\endgroup
      \else \csarg\ifx{LaLaEnd\CurrentComment Test}\test
            \edef\html@next{\endgroup\noexpand\end{\CurrentComment}}
      \else \let\html@next\xComment
      \fi \fi \html@next}
}
\makeatother

\def\includecomment
 #1{\expandafter\def\csname#1\endcsname{}%
    \expandafter\def\csname end#1\endcsname{}}
\def\excludecomment
 #1{\expandafter\def\csname#1\endcsname{\ThrowAwayComment{#1}}%
    {\escapechar=-1\relax
     \csarg\xdef{PlainEnd#1Test}{\string\\end#1}%
     \csarg\xdef{LaLaEnd#1Test}{\string\\end\string\{#1\string\}}%
    }}

%  Define environments that ignore their contents.
\excludecomment{comment}
\excludecomment{rawhtml}
\excludecomment{htmlonly}

%  Hypertext commands etc. This is a condensed version of the html.sty
%  file supplied with LaTeX2HTML by: Nikos Drakos <nikos@cbl.leeds.ac.uk> &
%  Jelle van Zeijl <jvzeijl@isou17.estec.esa.nl>. The LaTeX2HTML documentation
%  should be consulted about all commands (and the environments defined above)
%  except \xref and \xlabel which are Starlink specific.

\newcommand{\htmladdnormallinkfoot}[2]{#1\footnote{#2}}
\newcommand{\htmladdnormallink}[2]{#1}
\newcommand{\htmladdimg}[1]{}
\newcommand{\hyperref}[4]{#2\ref{#4}#3}
\newcommand{\htmlref}[2]{#1}
\newcommand{\htmlimage}[1]{}
\newcommand{\htmladdtonavigation}[1]{}

\newenvironment{latexonly}{}{}
\newcommand{\latex}[1]{#1}
\newcommand{\html}[1]{}
\newcommand{\latexhtml}[2]{#1}
\newcommand{\HTMLcode}[2][]{}

%  Starlink cross-references and labels.
\newcommand{\xref}[3]{#1}
\newcommand{\xlabel}[1]{}

%  LaTeX2HTML symbol.
\newcommand{\latextohtml}{\LaTeX2\texttt{HTML}}

%  Define command to re-centre underscore for Latex and leave as normal
%  for HTML (severe problems with \_ in tabbing environments and \_\_
%  generally otherwise).
\renewcommand{\_}{\texttt{\symbol{95}}}

% -----------------------------------------------------------------------------
%  Debugging.
%  =========
%  Remove % on the following to debug links in the HTML version using Latex.

% \newcommand{\hotlink}[2]{\fbox{\begin{tabular}[t]{@{}c@{}}#1\\\hline{\footnotesize #2}\end{tabular}}}
% \renewcommand{\htmladdnormallinkfoot}[2]{\hotlink{#1}{#2}}
% \renewcommand{\htmladdnormallink}[2]{\hotlink{#1}{#2}}
% \renewcommand{\hyperref}[4]{\hotlink{#1}{\S\ref{#4}}}
% \renewcommand{\htmlref}[2]{\hotlink{#1}{\S\ref{#2}}}
% \renewcommand{\xref}[3]{\hotlink{#1}{#2 -- #3}}
%end{latexonly}
% -----------------------------------------------------------------------------
% ? Document specific \newcommand or \newenvironment commands.

% degrees symbol
\newcommand{\dgs}{\hbox{$^\circ$}} 
\begin{htmlonly}
\newcommand{\dgs}{{\rawhtml &deg;}} 
\end{htmlonly}

% arcminute symbol
\newcommand{\arcm}{\hbox{$^\prime$}} 
\begin{htmlonly}
\newcommand{\arcm}{{\rawhtml &acute;}} 
\end{htmlonly}

% arcsec symbol
\newcommand{\arcsec}{\arcm\hskip -0.1em\arcm}
\begin{htmlonly}
\newcommand{\arcsec}{{\rawhtml &quot;}} 
\end{htmlonly}

% decimal-degree symbol
\newcommand{\udeg}{\hskip-0.3em\dgs\hskip-0.08em}
\begin{htmlonly}
\newcommand{\udeg}{{\rawhtml &deg;}} 
\end{htmlonly}

% decimal-arcsecond symbol
\newcommand{\uarcs}{\hskip-0.27em\arcsec\hskip-0.02em}  
\begin{htmlonly}
\newcommand{\uarcs}{{\rawhtml &quot;}}
\end{htmlonly}

\newcommand{\CCDPACK}{{\footnotesize CCDPACK}}
\newcommand{\CONVERT}{{\footnotesize CONVERT}}
\newcommand{\CURSA}{{\footnotesize CURSA}}
\newcommand{\GAIA}{{\footnotesize GAIA}}
\newcommand{\KAPPA}{{\footnotesize KAPPA}}
\newcommand{\ORACDR}{{\footnotesize ORAC-DR}}
\newcommand{\PHOTOM}{{\footnotesize PHOTOM}}
\newcommand{\PISA}{{\footnotesize PISA}}
\newcommand{\FITSref}{\htmladdnormallink{FITS}{http://fits.gsfc.nasa.gov/}}

% +
%  Name:
%     SST.TEX

%  Purpose:
%     Define LaTeX commands for laying out Starlink routine descriptions.

%  Language:
%     LaTeX

%  Type of Module:
%     LaTeX data file.

%  Description:
%     This file defines LaTeX commands which allow routine documentation
%     produced by the SST application PROLAT to be processed by LaTeX and
%     by LaTeX2html. The contents of this file should be included in the
%     source prior to any statements that make of the sst commnds.

%  Notes:
%     The commands defined in the style file html.sty provided with LaTeX2html
%     are used. These should either be made available by using the appropriate
%     sun.tex (with hypertext extensions) or by putting the file html.sty
%     on your TEXINPUTS path (and including the name as part of the
%     documentstyle declaration).

%  Authors:
%     RFWS: R.F. Warren-Smith (STARLINK)
%     PDRAPER: P.W. Draper (Starlink - Durham University)

%  History:
%     10-SEP-1990 (RFWS):
%        Original version.
%     10-SEP-1990 (RFWS):
%        Added the implementation status section.
%     12-SEP-1990 (RFWS):
%        Added support for the usage section and adjusted various spacings.
%     8-DEC-1994 (PDRAPER):
%        Added support for simplified formatting using LaTeX2html.
%     {enter_further_changes_here}

%  Bugs:
%     {note_any_bugs_here}

% -

%  Define length variables.
\newlength{\sstbannerlength}
\newlength{\sstcaptionlength}
\newlength{\sstexampleslength}
\newlength{\sstexampleswidth}

%  Define a \tt font of the required size.
\latex{\newfont{\ssttt}{cmtt10 scaled 1095}}
\html{\newcommand{\ssttt}{\tt}}

%  Define a command to produce a routine header, including its name,
%  a purpose description and the rest of the routine's documentation.
\newcommand{\sstroutine}[3]{
   \goodbreak
   \rule{\textwidth}{0.5mm}
   \vspace{-7ex}
   \newline
   \settowidth{\sstbannerlength}{{\Large {\bf #1}}}
   \setlength{\sstcaptionlength}{\textwidth}
   \setlength{\sstexampleslength}{\textwidth}
   \addtolength{\sstbannerlength}{0.5em}
% Changed to -1.0 from -2.0 because there is only space for one banner.
   \addtolength{\sstcaptionlength}{-1.0\sstbannerlength}
% Changed to -2.5 from -5.0 because there is only space for one banner.
   \addtolength{\sstcaptionlength}{-2.5pt}
   \settowidth{\sstexampleswidth}{{\bf Examples:}}
   \addtolength{\sstexampleslength}{-\sstexampleswidth}
   \parbox[t]{\sstbannerlength}{\flushleft{\Large {\bf #1}}}
   \parbox[t]{\sstcaptionlength}{\center{\Large #2}}
%   \parbox[t]{\sstbannerlength}{\flushright{\Large {\bf #1}}}
   \begin{description}
      #3
   \end{description}
}

%  Format the description section.
\newcommand{\sstdescription}[1]{\item[Description:] #1}

%  Format the usage section.
\newcommand{\sstusage}[1]{\pagebreak[3] \item[Usage:] \mbox{} \\[1.3ex] {\ssttt #1}}

%  Format the invocation section.
\newcommand{\sstinvocation}[1]{\item[Invocation:]\hspace{0.4em}{\tt #1}}

%  Format the arguments section.
\newcommand{\sstarguments}[1]{
   \item[Arguments:] \mbox{} \\
   \vspace{-3.5ex}
   \begin{description}
      #1
   \end{description}
}

%  Format the returned value section (for a function).
\newcommand{\sstreturnedvalue}[1]{
   \item[Returned Value:] \mbox{} \\
   \vspace{-3.5ex}
   \begin{description}
      #1
   \end{description}
}

%  Format the parameters section (for a ORAC-DR recipe).
\newcommand{\sstparameters}[1]{
   \goodbreak
   \item[Configurable Steering Parameters:] \mbox{} \\
   \vspace{-3.5ex}
   \begin{description}
      #1
   \end{description}
}

%  Format the examples section.
\newcommand{\sstexamples}[1]{
   \goodbreak
   \item[Examples:] \mbox{} \\
   \vspace{-3.5ex}
   \begin{description}
      #1
   \end{description}
}

%  Define the format of a subsection in a normal section.
\newcommand{\sstsubsection}[1]{ \item[{#1}] \mbox{} \\}

%  Define the format of a subsection in the examples section.
\newcommand{\sstexamplesubsection}[2]{\sloppy \item[\parbox{\sstexampleslength}{\ssttt #1}] \mbox{} \\ #2 }

%  Format the notes section.
\newcommand{\sstnotes}[1]{\pagebreak[3] \item[Notes:] \mbox{} \\[1.3ex] #1}

%  Provide a general-purpose format for additional (DIY) sections.
\newcommand{\sstdiytopic}[2]{\goodbreak \item[{\hspace{-0.35em}#1\hspace{-0.35em}:}] \mbox{} \\[1.3ex] #2}

%  Format the implementation status section.
\newcommand{\sstimplementationstatus}[1]{
   \pagebreak[3] \item[{Implementation Status:}] \mbox{} \\[1.3ex] #1}

%  Format the bugs section.
\newcommand{\sstbugs}[1]{\item[Bugs:] #1}

%  Specify a variant of the itemize environment where the top separation
%  is reduced.  It is needed because a \vspace is ignored in the
%  \sstitemlist command.
\newenvironment{sstitemize}{%
  \vspace{-4.3ex}\begin{itemize}}{\end{itemize}}

%  Format a list of items while in paragraph mode.
\newcommand{\sstitemlist}[1]{
  \mbox{} \\
  \vspace{-3.5ex}
  \begin{sstitemize}
     #1
  \end{sstitemize}
}

%  Format a list of items while in paragraph mode, and where there
%  is a heading, thus the negative vertical space is not needed.
\newcommand{\ssthitemlist}[1]{
  \mbox{} \\
  \vspace{-3.5ex}
  \begin{itemize}
     #1
  \end{itemize}
}

%  Define the format of an item.
\newcommand{\sstitem}{\item}

%  Now define html equivalents of those already set. These are used by
%  latex2html and are defined in the html.sty files.
\begin{htmlonly}

%  sstroutine.
   \newcommand{\sstroutine}[3]{
      \subsection{#1\xlabel{#1}-\label{#1}#2}
      \begin{description}
         #3
      \end{description}
   }

%  sstdescription
   \newcommand{\sstdescription}[1]{\item[Description:]
      \begin{description}
         #1
      \end{description}
      \\
   }

%  sstusage
   \newcommand{\sstusage}[1]{\item[Usage:]
      \begin{description}
         {\ssttt #1}
      \end{description}
      \\
   }

%  sstinvocation
   \newcommand{\sstinvocation}[1]{\item[Invocation:]
      \begin{description}
         {\ssttt #1}
      \end{description}
      \\
   }

%  sstarguments
   \newcommand{\sstarguments}[1]{
      \item[Arguments:] \\
      \begin{description}
         #1
      \end{description}
      \\
   }

%  sstreturnedvalue
   \newcommand{\sstreturnedvalue}[1]{
      \item[Returned Value:] \\
      \begin{description}
         #1
      \end{description}
      \\
   }

%  sstparameters
   \newcommand{\sstparameters}[1]{
      \item[Configurable Steering Parameters:] \\
      \begin{description}
         #1
      \end{description}
      \\
   }

%  sstexamples
   \newcommand{\sstexamples}[1]{
      \item[Examples:] \\
      \begin{description}
         #1
      \end{description}
      \\
   }

%  sstsubsection
   \newcommand{\sstsubsection}[1]{\item[{#1}]}

%  sstexamplesubsection
   \newcommand{\sstexamplesubsection}[2]{\item[{\ssttt #1}] #2}

%  sstnotes
   \newcommand{\sstnotes}[1]{\item[Notes:]
      \begin{description}
         #1
      \end{description}
      \\
   }

%  sstdiytopic
   \newcommand{\sstdiytopic}[2]{\item[{#1}:]
      \begin{description}
         #2
      \end{description}
      \\
   }

%  sstimplementationstatus
   \newcommand{\sstimplementationstatus}[1]{\item[Implementation Status:]
      \begin{description}
         #1
      \end{description}
      \\
   }

%  sstitemlist
   \newcommand{\sstitemlist}[1]{
      \begin{itemize}
         #1
      \end{itemize}
      \\
   }

%  ssthitemlist
   \newcommand{\ssthitemlist}[1]{
      \begin{itemize}
         #1
      \end{itemize}
      \\
   }
\end{htmlonly}

%  End of "sst.tex" layout definitions.

% ? End of document specific commands
% -----------------------------------------------------------------------------
%  Title Page.
%  ===========
\renewcommand{\thepage}{\roman{page}}
\begin{document}
\thispagestyle{empty}

%  Latex document header.
%  ======================
\begin{latexonly}
   CCLRC / \textsc{Rutherford Appleton Laboratory} \hfill \textbf{\stardocname}\\
   {\large Particle Physics \& Astronomy Research Council}\\
   {\large Starlink Project\\}
   {\large \stardoccategory\ \stardocnumber}
   \begin{flushright}
   \stardocauthors\\
   \stardocdate
   \end{flushright}
   \vspace{-4mm}
   \rule{\textwidth}{0.5mm}
   \vspace{5mm}
   \begin{center}
   {\Huge\textbf{\stardoctitle \\ [2.5ex]}}
   {\LARGE\textbf{\stardocversion \\ [4ex]}}
   {\Huge\textbf{\stardocmanual}}
   \end{center}
   \vspace{5mm}

% ? Add picture here if required for the LaTeX version.
%   e.g. \includegraphics[scale=0.3]{filename.ps}
\begin{center}
\includegraphics[width=1.0in]{sun232_logo.eps}
\end{center}
% ? End of picture

% ? Heading for abstract if used.
   \vspace{10mm}
   \begin{center}
      {\Large\textbf{Abstract}}
   \end{center}
% ? End of heading for abstract.
\end{latexonly}

%  HTML documentation header.
%  ==========================
\begin{htmlonly}
   \xlabel{}
   \begin{rawhtml} <H1> \end{rawhtml}
      \stardoctitle\\
      \stardocversion\\
      \stardocmanual
   \begin{rawhtml} </H1> <HR> \end{rawhtml}

% ? Add picture here if required for the hypertext version.
\includegraphics[width=1.0in]{sun232_logo.eps}
%   e.g. \includegraphics[scale=0.7]{filename.ps}
% ? End of picture

   \begin{rawhtml} <P> <I> \end{rawhtml}
   \stardoccategory\ \stardocnumber \\
   \stardocauthors \\
   \stardocdate
   \begin{rawhtml} </I> </P> <H3> \end{rawhtml}
      \htmladdnormallink{CCLRC / Rutherford Appleton Laboratory}
                        {http://www.cclrc.ac.uk} \\
      \htmladdnormallink{Particle Physics \& Astronomy Research Council}
                        {http://www.pparc.ac.uk} \\
   \begin{rawhtml} </H3> <H2> \end{rawhtml}
      \htmladdnormallink{Starlink Project}{http://www.starlink.rl.ac.uk/}
   \begin{rawhtml} </H2> \end{rawhtml}
   \htmladdnormallink{\htmladdimg{source.gif} Retrieve hardcopy}
      {http://www.starlink.rl.ac.uk/cgi-bin/hcserver?\stardocsource}\\

%  HTML document table of contents. 
%  ================================
%  Add table of contents header and a navigation button to return to this 
%  point in the document (this should always go before the abstract \section). 
  \label{stardoccontents}
  \begin{rawhtml} 
    <HR>
    <H2>Contents</H2>
  \end{rawhtml}
  \htmladdtonavigation{\htmlref{\htmladdimg{contents_motif.gif}}
        {stardoccontents}}

% ? New section for abstract if used.
  \section{\xlabel{abstract}Abstract}
% ? End of new section for abstract
\end{htmlonly}

% -----------------------------------------------------------------------------
% ? Document Abstract. (if used)
%  ==================
\stardocabstract
% ? End of document abstract
% -----------------------------------------------------------------------------
% ? Latex document Table of Contents (if used).
%  ===========================================
  \newpage
  \begin{latexonly}
    \setlength{\parskip}{0mm}
    \tableofcontents
    \setlength{\parskip}{\medskipamount}
    \markboth{\stardocname}{\stardocname}
  \end{latexonly}
% ? End of Latex document table of contents
% -----------------------------------------------------------------------------
\cleardoublepage
\renewcommand{\thepage}{\arabic{page}}
\setcounter{page}{1}

\section{\xlabel{introduction}Introduction\label{introduction}}

\ORACDR\ is a data-reduction pipeline at UKIRT, part of the
\htmladdnormallink{ORAC system}{http://www.stsci.edu/stsci/meetings/adassVII/bridgera.html}
The pipeline reduces and displays multi-frame
observations soon after they are read from the detector.  This allows
observers to assess the quality and suitability of their data in near
real time.  Yet \ORACDR\ is capable of producing publication-quality
results.  

\ORACDR\ is suitable for `offline' data reduction at your home
institution too.  There are many reasons why you may wish to use
\ORACDR\ in this fashion.  For instance, you may have come back from
UKIRT with only the raw observations; or there was an error in a
telescope `exec' mixing the groups of observations; or some data were
reduced with a basic algorithm for speed at the telescope, and now you
want to do a more-careful job.  Hence \ORACDR\ is now available on
Starlink.

\xref{SUN/230}{sun230}{} presents an overview of \ORACDR,
general facilities like its display system, and it explains the
differences between a pipeline and a traditional reduction package.
Put briefly, \ORACDR\ uses a few data headers to direct the data
reduction.  Amongst these headers is the name of a {\em recipe}.  A
recipe is a series of high-level instructions such as ``make a
mosaic'' or ``divide by a flat'' that reduces an {\em observation\/}
comprising one or more data frames.  The implementation of each of
these instructions is through a Perl script---called a
primitive---which calls Starlink packages such as \CCDPACK\ and \KAPPA,
to actually do the processing of the bulk data.

This document describes how to use \ORACDR\ software on Starlink to
reduce data from the UKIRT imaging instruments: UFTI and IRCAM. It
outlines the various algorithms used in the recipes, and includes
detailed recipe documentation in the appendix.  Besides the standard
reduction recipes, this manual describes how you can customise recipes
to suit your preferences, and how to correct errors in the headers of
your data frames.

There is a complementary document \xref{SUN/231}{sun231}{} which
addresses the reduction of SCUBA data with \ORACDR.

\section{\xlabel{Using_the_pipeline}Using the
pipeline\label{Using_the_pipeline}}

\subsection{Setting up \ORACDR}

Before you can run the pipeline you have to tell \ORACDR\ for which
instrument you wish to reduce data, the observation date, and the
directory containing the raw data, and wher you want the processed
data to be written.  There are two options.

\begin{itemize}

\item  The first needs your data to conform to the directory-naming
convention of the instrument at UKIRT.  This will be the case if you
simply unpack the archive written by the {\bf uktape} utility.  In
this case enter

\begin{verbatim}
      % setenv ORAC_DATA_ROOT <root_data_directory>
      % oracdr_<instrument> <date>
\end{verbatim}
where {\tt $<$root\_data\_directory$>$} is the directory in which you
unpacked the data from the tape, {\tt $<$instrument$>$} is either {\tt
ufti} or {\tt ircam}, and {\tt$<$date$>$} is the UT date in the format
YYYYMMDD.  Note that the \texttt{\%} represent the UNIX shell's
prompt, which you do not type.  The commands must be entered in the
above order.  

For example, the standard location for raw UFTI data is
{\tt ufti\_data/YYYYMMDD/raw/}, and {\tt ufti\_data/YYYYMMDD/reduced/}
for the corresponding reduced data.  So if your data are stored in
{\tt /home/users/abc/data/UKIRT/ufti\_data/19991108/raw/} you should
enter the following.

\begin{verbatim}
      % setenv ORAC_DATA_ROOT  /home/users/abc/data/UKIRT/
      % oracdr_ufti 19991108
\end{verbatim}
to enable the pipeline for UFTI data taken on 1999 November 8.

\item The second option is where your raw and reduced data are to be
in arbitrary directories.  Type the following

\begin{verbatim}
      % oracdr_<instrument> <date>
      % setenv ORAC_DATA_IN <raw_data_directory>
      % setenv ORAC_DATA_OUT <reduced_data_directory>
\end{verbatim}

Here is an example for IRCAM data.
\begin{verbatim}
      % oracdr_ircam 19990328
      % setenv ORAC_DATA_IN /export/data/mjc/asteroid/night1
      % setenv ORAC_DATA_OUT /home/scratch/mjc/reduced
\end{verbatim}

\end{itemize}
In the first case {\tt \$ORAC\_DATA\_IN} and {\tt \$ORAC\_DATA\_OUT} are
still defined, but in terms of the root directory.  For instance, re-using
the earlier example with UFTI for UT date 1999 November 8,
{\tt \$ORAC\_DATA\_IN} points to {\tt \$ORAC\_DATA\_ROOT/ufti\_data/19991108/raw/}.

\ORACDR\ operates in {\tt \$ORAC\_DATA\_OUT}, irrespective of
what your current directory is when you invoke it.  Raw data are
either copied (IRCAM) to {\tt \$ORAC\_DATA\_OUT}, or raw \FITSref\
files are converted to NDFs (UFTI) in this directory.  Your current
directory remains unchanged.

It is highly recommended to work in directories on discs local to
the computer running the pipeline.  Processing over NFS-served drives
can many times slower and degrades the performance seen by other users.

\subsection{\xlabel{running_the_pipeline}Running the
pipeline\label{running_the_pipeline}}

To run the pipeline, you use the {\bf oracdr} command.  This has a
number of qualifiers described fully in \xref{SUN/230}{sun230}{oracdr}.
There is online help too; enter

\begin{verbatim}
      % oracdr -h
\end{verbatim}
for a list of the options.

Unlike using \ORACDR\ at UKIRT, you are unlikely to need the looping ({\tt -loop}
option) for offline processing, as all the data exist.  Thus the most
important qualifiers are {\tt -list} and {\tt -from}, which specify
the frames to process; and the recipe name.

\begin{verbatim}
      % oracdr -from 42
\end{verbatim}
will process frames f19991108\_00042 until the end of the night's data
(assuming the earlier {\bf orac\_ufti} command), running the recipes
given by each frame's header (RECIPE keyword).  More likely is that
you provide a list of selected observations.  The following example

\begin{verbatim}
      % oracdr -list 41:49,51:59 JITTER9_SELF_FLAT
\end{verbatim}
processes frames from 41 to 49 inclusive and 51 to 59 inclusive,
invoking the {\tt JITTER9\_SELF\_FLAT} recipe, and overriding
the RECIPE header.

\begin{verbatim}
      % oracdr -list 5:7,23,33
\end{verbatim}
woild reduce the frames 5, 6, 7, 23, and 33.  This is most likely
to be applicable to a series of dark frames.

There is a hazard with the {\tt -list} option.  Take care to select a
complete set of frames associated with an observation.  A common error
is to include a dark frame accidently.  Check the log in the raw data
directory; it has file extension {\tt .nightlog}.  If you do not have
a log, it is easy to create one.

\begin{verbatim}
      % oracdr -noeng -from 1 -skip NIGHT_LOG
\end{verbatim}
This will create a log called {\tt \$ORAC\_DATA\_IN/$<$date$>$.nightlog}
for the current UT date.

\subsection{\xlabel{display}Display\label{display}}

\ORACDR\ optionally lets you inspect the raw frames, and the processed 
data as they are created happens.  There is a variety 
of graphics methods available, including histograms and contour plots,
if you choose a \KAPPA\ GWM widget.  Most people prefer a simple scaled
image display with \GAIA.  The selection of frame types to display, 
where they should appear, and how they are scaled are configurable
using a simple text file or a special GUI tool {\bf oracdisp}.  See
\xref{SUN/230}{sun230}{display_system} for details and examples.

Processing offline, there is less need to see the data displayed in real time.
If you wish to accelerate the processing switch off the display option.

\begin{verbatim}
      % oracdr -nodisplay ...
\end{verbatim}

If you do want to display a recommendation is to create two GAIA windows
displaying images using autoscaled limits.  This first could be for
raw and flat-fielded data, and the second for the mosaics.  You are
likely to want to interact with the latter, using \GAIA's toolboxes.
Your {\tt \$ORAC\_DATA\_OUT/disp.dat} could look like this.

\begin{verbatim}
      # Send raw frame to first GAIA window 
      num type=image tool=gaia region=0 window=0 autoscale=1 zautoscale=1

      # Send flatfielded frame to first GAIA window.
      ff  type=image tool=gaia region=0 window=0 autoscale=1 zautoscale=1

      # Send mosaic frame to second GAIA window.
      mos type=image tool=gaia region=0 window=1 autoscale=1 zautoscale=1
\end{verbatim}

\section{\xlabel{features_of_the_primitives}Features of the
Primitives\label{features_of_the_primitives}}

Primitives are the Perl scripts which actually call the applications
to do most of the data processing. All but two of the imaging recipes are
independent of the instrument.  Appropriate header information is
extracted for IRCAM and UFTI.

\subsection{\xlabel{preliminaries}Preliminaries\label{preliminaries}}

There a few operations applied to all frames.  The first sets the
origin of the frame so that frame pixels retain the detector pixel
indices.  It then becomes possible to use a full-sized bad-pixel mask
or flat field on any subset of a detector's pixel grid.

The second is to switch on history recording.  It is recommended to
leave this enabled, since it provides a record of the processing steps
of your final mosaics.  Otherwise the pipeline becomes sometime of 
a black box.  Use the \KAPPA\ command
\xref{{\bf hislist}}{sun95}{HISLIST} to list the history records.

IRCAM has two additional steps.  The pipeline has to copy the files
because the original files are normally write protected.  It finally
corrects for detector non-linearity.

\subsection{\xlabel{flat_creation}Flat creation\label{flat_creation}}

Frames are optionally cleaned to remove extreme outliers ($\pm$3 or
6$\sigma$ about the mean in 15$\times$15-pixel neighbourhood) iterated
three times.  The data are then normalised, combined pixel by pixel
using a median, and the combined array normalised to a have a mean of
one.  It is possible to use another statistic for the combination,
such as a clipped median.

\subsection{\xlabel{object_masking}Object masking\label{object_masking}}

In recipes which make a flat using the frames taken of the targets,
the so-called self flat, any sources present can bias the flat field, and
result in blotchy mosaics.  The full versions of such recipes, as
opposed to the \_BASIC versions, and SKY\_MASKED greatly reduce these
artifacts using the following algorithm.  After the application of the
approximate self-flat field, an inventory is made of objects having at
least 12 connected pixels above one sigma above sky.  The locations,
shapes, orientations and sizes are used to make a mask.  The mask is
applied to the dark-subtracted frames and a new flat created.  As
the outer parts of bright objects often leave residual unmasked
blobs, a circular central occulting mask is used.  The diameter is
normally 7 arcseconds, but it can be adjusted through the OCCULT argument
of primitive \_MAKE\_OBJECTS\_MASK\_.  In the QUADRANT\_JITTER recipe
the central mask's diameter equals the length of the shorter side
of a quadrant.

The improved flat typically shows a uniformity at $\sim$0.02\% of the
sky. It is this flat which produces the flat-field frames for
mosaicking.  Systematic errors in the sky---a major uncertainty in
infra-red point-source photometry---are also reduced significantly by
this algorithm.

\subsection{\xlabel{automatic_registration}Automatic
Registration\label{automatic_registration}}

This makes an inventory of the sources above a threshold in each
frame.  It then performs a \xref{pattern
recognition}{sun139}{mosaicing} to identify common features in
jittered frames.  If the fraction of common objects is under 40\% or
the total is fewer than three, the registration fails, and so the
script resorts to reading the offsets stored in the \FITSref\ headers, or
matching a central bright object in certain recipes.  Using telescope
offsets can lead to trailed sources, as occurred with the
\xref{{\footnotesize IRCAMDR} package}{sun41}{}.  The improved
registration leads to the detection of fainter sources and more-accurate
measurement thereof.

\subsection{\xlabel{mosaicking}Mosaicking\label{mosaicking}}

This is fully automated.  No longer do you have to measure
painstakingly centroids and manually tile to form mosaics from the
jittered frames.  The jitter offsets are sufficiently small to permit
shift of Cartesian co-ordinates to register.  The offsets are derived
from the automatic registration.  No correction for the detector
orientation is applied, since it degrades the quality of the data,
despite the small rotation 
\begin{latexonly}
($\sim$1$\dgs$)
\end{latexonly}
\begin{htmlonly}
($\sim$1 degree)
\end{htmlonly}
normally involved.  Allowance for this misalignment with the cardinal
directions will be through the \htmladdnormallink{FITS world co-ordinate
system}{http://fits.gsfc.nasa.gov/documents.html#WCS}. In the meantime
should you need the rotation angle, pipe the output of KAPPA's
\xref{{\bf fitslist}}{sun95}{FITSLIST} function into grep.

\begin{verbatim}
      % fitslist <your_frame> | grep CROTA2
\end{verbatim}

The flat-fielded frame can either be resampled to give sub-pixel
registration, or to the nearest pixel. The latter is much faster and
is adequate for most uses of UFTI and IRCAM/TUFTI.  For older
undersampled IRCAM3 data
\begin{latexonly}
(0.$\uarcs$286
\end{latexonly}
\begin{htmlonly}
(0.286 arcseconds
\end{htmlonly}
per pixel) resampling is recommended.

The mosaicking uses the \xref{\CCDPACK\ algorithm}{sun139}{mosaicing}
of its \xref{{\bf makemos}}{sun139}{MAKEMOS} command.  Only zero-point
shifts of intensity are applied to the resampled frames to create the
mosaic.  For most cases the comparison is of the sky levels as sky
pixels dominate.  This comparison is repeated for all pairs.  {\bf
makemos} then finds the most mutually consistent set of additive
corrections, weighting appropriately, to make the smoothest mosaic
given the data. The first frame, which for the UFTI execs is the
central (offset 0 0) frame, has no additive correction applied.  The
mosaic generation adjusts the zero-point of the jittered frames.
Another way of looking at it is that mosaicking attempts to remove the
sky variations.  The additive corrections are normally quite small,
like a few counts, but over longer integrations or in the thermal
regime they can amount to a few score.

There is no normalisation to counts per second in the mosaic. The
mosaic's signal corresponds to that of the first frame, thus the
exposure time of the mosaic is that of one individual frame.  The
recipes assume that you have used their corresponding `execs', and
hence have not changed the exposure time during a jitter.  The
exposure time (header EXP\_TIME) is propagated from the first frame to
the mosaic.  Where multiple frames combined to create a mosaic pixel
the signal-to-noise ratio corresponds to the combined integration
time.  The integration time (keyword INT\_TIME) is the number of
coadds times the exposure time per coadd.

Depending on the recipe, the mosaic may be trimmed to the dimensions
of a raw frame.  Mosaicking removes virtually all the bad pixels for
standard stars where the jitter offsets are small.

A mosaic forms for each cycle of the recipe, e.g. all nine frames in
a JITTER9\_SELF\_FLAT.  For multiple cycles, an integrated `grand'
mosaic forms of improving signal-to-noise.  To avoid the build up of
bad pixels from cosmic rays, bad pixels are interpolated before the
addition.  This may result in some strange stripes in the top-left
corner of UFTI frames where no interpolation can occur.  Those pixels
are bad in all frames and should be ignored.  The EXP\_TIME for the
integrating mosaic is the sum of the exposure times of the
contributing mosaics.  Again the signal is not divided by the exposure
time.

\subsection{\xlabel{aperture_photometry}Aperture
Photometry}\label{aperture_photometry}

The recipes with an \_APHOT suffix perform aperture photometry on
the mosaic and the contributing flat-fielded frames.  The method
assumes that the target, usually a standard star, is approximately
centrally located after allowing for the jitter offsets.  If you have
data where the star lies outside the aperture, it is possible to
apply an offset.  See the XOFF and YOFF arguments of primitive
\_FIXED\_APERTURE\_PHOTOMETRY\_ in \_APHOT\_MAG\_ to adjust the
aperture's position.

The photometry is through a circular aperture located at the centroid
of the source, with the sky measured from a concentric annulus outside
the aperture.  The typical aperture size is 5 or 6 arcseconds,
depending on the recipe.  The annulus is typically 9 to 15 arcseconds
diameter.  The default estimator of the sky flux is the mode.  The
photometry accounts for fractional pixels at the aperture edge but
without allowance for the local gradient.

The magnitudes are given by the expression $-2.5 * $log10(abs(counts)
per second exposure time).  Therefore negative sources can be
measured too, as presented by the NOD recipes.  The photometry also
yields an internal error determined from the sky variance. 

A case- and space-insensitive comparison of the object name with the
entries in a table provides a catalogue magnitude in $J$, $H$, or $K$
for a standard star.  Also a mean extinction is applied for the mean
of the start and end airmasses.  Thus the primitive calculates an
approximate zero point.  For accurate photometry the actual extinction
coefficients should be determined.  As the output from the photometry
is a \xref{small text list}{sun190}{STLREF}, you can use
\xref{{\bf catphotomfit} command}{sun190}{PHOTCAL} of the \CURSA\ package
to achieve this.

\subsection{\xlabel{tidying}Tidying\label{tidying}}

Each recipe has a tidy procedure, which removes unnecessary
intermediate frames when the recipe no longer requireds them. Retained
are the raw data, flat-fielded frames, and the mosaics.  Most of the
intermediate small text files are removed in individual primitives,
but some do persist.

\section{\xlabel{customising_recipes}Customising
Recipes\label{customising_recipes}}

At the time of writing, the {\em \ORACDR\ Programmer's Guide\/} User
Note is unavailable.  For most purposes observers wishing to make or
modify existing scripts can get by without this document.

\subsection{\xlabel{search_paths}Search paths\label{search_paths}}

\ORACDR\ allows you to create your own recipes and primitives, or
modify those provided as part of the package.  In either case you must
tell \ORACDR\ where your recipes and/or primitives are stored.  This
is achieved through two environment variables.  {\tt\$ORAC\_RECIPE\_DIR}
should equate to the directory containing your recipes.
{\tt \$ORAC\_PRIMITIVE\_DIR} specifies the directory containing your
primitives.  Here's an example.

\begin{verbatim}
      % setenv ORAC_RECIPE_DIR /home/user/drmoan/recipes
      % setenv ORAC_PRIMITIVE_DIR /home/user/drmoan/primitives
\end{verbatim}

Once these environment variables are defined, \ORACDR\ first looks in
{\tt \$ORAC\_RECIPE\_DIR} or {\tt \$ORAC\_PRIMITIVE\_DIR} to find a recipe or
primitive respectively.  If the script is absent, \ORACDR\ looks in
the standard {\tt \$ORAC\_DIR} directories.

\subsection{\xlabel{anatomy_of_an_imaging_recipe}Anatomy of an imaging recipe}

There are documentation modules---a Starlink style between \verb/#+/ and
\verb/#-/ delimiters at the head, and a Perl POD (Plain Old
Documentation) at the foot.  Between these is the code.  This consists
of calls to primitives, sometimes with arguments.  Primitives have
uppercase names preceded and terminated by underscores, such as
\_DIVIDE\_BY\_FLAT\_.

\subsubsection{\xlabel{hello_primitives}Hello primitives}

The first of these primitives is either \_UFTI\_HELLO\_ or \_IRCAM\_HELLO\_.
They contain instrument-specific code and initialisation.  Those are
best left alone.  See 
\begin{htmlonly}
\htmlref{preliminaries}{preliminaries}
\end{htmlonly}
\begin{latexonly}
Section~\ref{preliminaries}
\end{latexonly}
for a description of what these primitives do.

Second there is a recipe-specific primitive such as
\_JITTER\_SELF\_FLAT\_HELLO\_.  This sets up \CCDPACK.  Two things you
might wish to change in the recipe hello script are listed below.

\begin{itemize}

\item Change the extent of the images.  If there is an instrumental defect
in some peripheral rows you might not want to use the full bounds
as given by headers RDOUT\_X1 ($x$ lower bound), RDOUT\_X2
($x$ upper bound), and RDOUT\_Y1 and RDOUT\_Y2 for $y$.  Suppose you
wanted to trim off the top three rows you could change the line.

\begin{verbatim}
    my $y2 = $Frm->hdr( "RDOUT_Y2" );
\end{verbatim}
to
\begin{verbatim}
    my $y2 = $Frm->hdr( "RDOUT_Y2" ) - 3;
\end{verbatim}

\item Switch on error propagation.  To save time at the telescope, the pipeline
does not keep track of the errors per pixel.  If you wish to know
realistic errors for your data, in the line defining the hidden
variables for the \CCDPACK\ setup, change the parameter GENVAR's
value from {\tt{no}} to {\tt{yes}}.

\end{itemize}

\subsubsection{\xlabel{steering_primitives}Steering primitive}

Third comes a `steering' primitive.  These are best left well alone.
They control when the various operations are performed.  What you can
safely adjust are the configurable steering parameters listed in the
recipe documentation.  In the main these parameters set the number of
frames processed in a pass through the recipe. 

Suppose that for some reason an observation of a nine-point jitter
self flat was aborted after seven positions.  If you try the recipe
stored in the headers some processing will occur, but it will not
include mosaic creation.  The final steps inclusing mosaic creation
occur once all nine frames are dark subtracted.  Now there are no
seven-point recipes to substitute on the command line.  You could make
your own seven-point recipe to reduce those data.  First make a new
recipe by copying the standard one.  Recipes are stored in
{\tt\$ORAC\_DIR/recipes/$<$instrument$>$}, where {\tt$<$instrument$>$} is
either {\tt UFTI} or {\tt IRCAM}.

\begin{verbatim}
    cd $ORAC_RECIPE_DIR
    cp $ORAC_DIR/recipes/UFTI/JITTER9_SELF_FLAT ./JITTER7_SELF_FLAT
\end{verbatim}
Next edit JITTER7\_SELF\_FLAT and alter the line

\begin{verbatim}
      _JITTER_SELF_FLAT_STEER_ NUMBER=9
\end{verbatim}
substituting {\tt 7} for the {\tt 9}.  The recipe will then generate
the self flat, flat field and make the mosaic once the seventh frame is
dark-subtracted.

\subsubsection{\xlabel{recipe_primitives}Recipe primitives}

After the steering primitive we come to the recipe-specific scripts
that actually perform the recipe.  The most likely and easiest things
you would change are to add arguments or modify argument values of the
primitives in the recipes.  For instance, you might wish to change the
aperture diameter for the aperture photometry.  To alter to a
10-arcsecond aperture, change the APERTURE argument's value of the
\_APHOT\_MAG\_ primitive like below.

\begin{verbatim}
      _APHOT_MAG_ APERTURE=10
\end{verbatim}

To obtain details of a primitive's arguments, use the {\bf oracman}
command.  Thus

\begin{verbatim}
      % oracman _MAKE_MOSAIC_NO_BOUND_
\end{verbatim}
will display the documentation for primitive
\_MAKE\_MOSAIC\_NO\_BOUND\_.  Space does not permit inclusion of the
documentation of the many primitives in this manual.  The sources are
stored in {\tt \$ORAC\_DIR/primitives/$<$instrument$>$}, where
{\tt$<$instrument$>$} is either {\tt UFTI} or {\tt IRCAM}.

While the simplest primitives just invoke a Starlink task, and updates
just that and are amenable to customisation, some are quite complex
especially for the registration.  They may invoke other primitives,
manipulate parameters and small data files so that the various tasks
connect to cope with a variety of circumstances.  The most likely
change you will want to make is to change the parameter values of a
Starlink task.  Armed with the reference documentation for the
application, say with a {\tt findme $<$application$>$}, it is easy
to change values or append further parameters.  

Here is an example.  Let us suppose you wanted to combine frames to
make a self flat, not with the median, since you have heard that a mean
trimmed of the most-extreme tenth of the values gives better results.
First copy \_MAKE\_FLAT\_FROM\_GROUP\_ to your primitives directory.

\begin{verbatim}
      cp $ORAC_DIR/primitives/UFTI/_MAKE_FLAT_FROM_GROUP_ $ORAC_PRIMITIVE_DIR
\end{verbatim}
Using an editor, find the first line in your copy of
\_MAKE\_FLAT\_FROM\_GROUP\_ commencing {\tt \$hidden}.  It should be
as follows.

\begin{verbatim}
      $hidden = "method=median sigmas=2.0 reset accept"; 
\end{verbatim}
Change this to

\begin{verbatim}
      $hidden = "method=trimmed alpha=0.1 sigmas=2.0 reset accept"; 
\end{verbatim}
to effect the change of statistic.  There is in fact a second line
assigning variable {\tt \$hidden} depending on argument CLEAN, and you
should make the same alteration there too.

If there is demand, additional arguments could be provided for
primitives, to simplify control.  Please contact the author if you
have suggestions for arguments and new recipes, or need help
customising your \ORACDR\ scripts.

\subsection{\xlabel{index_files}Index files\label{index_files}}

Once the pipeline has run for a bit you will find text files in {\tt
\$ORAC\_DATA\_OUT} called {\tt index.flat}, {\tt index.dark} amongst
others.  These list the calibration frames.  \ORACDR\ uses these to
find the most-recent, appropriate calibration.  For example, a flat
requires that the filter of the flat matches that of the frame being
flat fielded, and a dark must have the same exposure time as the
target frame; and both must have been taken in the same instrument
mode.

Here is an example of a flat index.

\begin{verbatim}
#FILTER MODE ORACTIME RDOUT_X1 RDOUT_X2 RDOUT_Y1 RDOUT_Y2
flat_23 Lp98 STARE 13.3484 1 256 1 256
flat_436 J98 flush_read 1.40859E+01 1 1024 1 1024
flat_78 K98 flush_read 7.23825E+00 1 1024 1 1024
flat_88 J98 flush_read 7.46111E+00 1 1024 1 1024
flat_93 H98 flush_read 7.58330E+00 1 1024 1 1024
flat_98 K98 flush_read 7.70559E+00 1 1024 1 1024
flat_133_subgrp H98 flush_read 8.99034E+00 1 1024 1 1024
flat_138_subgrp H98 flush_read 9.07139E+00 1 1024 1 1024
flat_138_subgrp_cycle1 H98 flush_read 9.12074E+00 1 1024 1 1024
flat_290_row0 K98 flush_read 1.13094E+01 1 1024 1 1024
flat_290_row1 K98 flush_read 1.15080E+01 1 1024 1 1024
\end{verbatim}
The first line contains the column headings.  ORACTIME is the UT in
decimal hours.

In general you should not manipulate these files.  Mis-editing can
lead to the calibration system breaking down.  If you must edit this
file, say to exclude a poor dark or an uneven flat, restrict yourself
deleting the line corresponding to that calibration file.  It's safer
to remove the calibration file and recreate a new one with the
calibration frames you want by running the pipeline.

If you want to nominate specific calibration frames, overriding those
selected from the calibration indices, there is a {\tt -calib} option
for the {\bf oracdr} command to do this.  See
\begin{latexonly}
the section on
\end{latexonly}
\xref{calibration options}{sun230}{calibration_options}
\begin{latexonly}
in SUN/230
\end{latexonly}
for examples.

\section{\xlabel{correcting_headers}Correcting headers}

There are reasons why you may need to edit some of the FITS headers
used by \ORACDR.

\begin{itemize}
\item  At the summit of Mauna Kea, it's easy to make mistakes.  One of 
the common ones is to make an error in the `exec'.  This can cause
frames to be in the wrong observation groups or be assigned the wrong
data-reduction recipe.

\item You may have made some trial observations before taking making
a longer integration through several cycles of a recipe.  Now you wish
to combine all the observations of a target to obtain the best
signal-to-noise.  
\end{itemize}

The main headers to change are
\begin{description}
\item [{\tt RECIPE}]---the data-reduction recipe;
\item [{\tt OBSNUM}]---the number of the frame, starting from 1 on
each night;
\item [{\tt GRPNUM}]---the group number, and should be
given by the frame number (OBSNUM) of its first member; and
\item [{\tt GRPMEM}]---the sequence number within the group, counting
from 1.
\end{description}
Three of them should obey the following relationship:
${\tt GRPNUM} = {\tt OBSNUM} - {\tt GRPMEM} + 1$.

For UFTI you can edit the raw FITS files.  {\tt \$ORAC\_DIR/bin/fitsmod.pl}
is a documented example Perl script to edit FITS headers.  The intention
is for you to make a copy and edit to suit your particular header-editing
requirements.

For IRCAM and UFTI, it's possible to edit the NDF's
\xref{FITS extension}{sun95}{se_fitsairlock} using \KAPPA's
\xref{{\bf fitsmod}}{sun95}{FITSMOD} command.  The command is a bit
long and the author regrets not defining a {\bf fitsupdate} synonym.
The following changes the GRPNUM keyword to have value 36 in the raw
NDF frame f19991108\_00042.

\begin{verbatim}
      % fitsmod f19991108_00042 grpnum update 36 \$C !
\end{verbatim}

It's possible to edit many files using a \xref{C-shell}{sc4}{} or Perl
script to edit a series of files very quickly.  If you do, it's better
to specify the values by keyword instead of position, like this

\begin{verbatim}
      % fitsmod ndf=f19991108_00042 edit=update keyword=grpnum value=36 \
                position=! comment=\$C
\end{verbatim}
because it is better insulated against change to {\bf fitsmod}.

\section{Acknowledgments}

\ORACDR\ was developed at the Joint Astronomy Centre by Frossie
Economou and Tim Jenness in collaboration with the UK Astronomy
Technology Centre as part of the ORAC project.  I should like to
thanks to members of the ORAC team, UKIRT staff and observers who made
suggestions for new or improved recipes.  Special thanks go to Gillian
Wright and Sandy Leggett for defining the initial specifications of
the UFTI scripts, and for subsequent discussions.  Thanks also to
Frossie Economou and Tim Jenness for answering my \ORACDR\ and Perl
questions, and for incorporating my requested enhancements into
\ORACDR.

\section{Copyright and License}

\ORACDR\ is copyright \copyright 1998--2000 PPARC (the UK Particle Physics
and Astronomy Research Council).  It is distributed by Starlink
under the GNU General Public License as published by the Free Software
Foundation.

Whenever you have used \ORACDR\ as part of a publication, please give
an acknowledgment to \ORACDR\ in the paper.  This will help us assess
the usage of \ORACDR.

\newpage
\appendix
\section{\xlabel{UFTI_file_suffices}UFTI file
suffices}\label{UFTI_file_suffices}%
\index{ufti_file_suffices}

\subsection*{Description}%
\index{Description}

Files generated during UFTI data reduction have suffices denoting the
processing step that created them.  This appendix contains a list with
short descriptions of what they mean.  Most will be removed once a
recipe has finished using them.  So you will probably only see these
files if you list the contents of directory {\tt \$ORAC\_DIR} while the
pipeline is running, or you interrupt the pipeline with {\tt CTRL/C},
or something has gone wrong with the the recipe and pipeline has
aborted.

\subsection*{UFTI Frame suffices}%
\index{UFTI Frame suffices}

\begin{tabular}{llp{90mm}}
\hline
  Suffix  & Stands for        & Description \\ \hline
{\tt\_bp} & Bad Pixel         & Co-added with the bad-pixel mask \\
{\tt\_cl} & CLone             & Modifiable copy of IRCAM raw data \\
{\tt\_db} & De-Biassed        & The bias is actually zero, but it sets
                                up various \xref{\CCDPACK}{sun139}{} things
                                for later processing \\
{\tt\_dg} & De-Glitched       & Bad pixels replaced by median of neighbours \\
{\tt\_dk} & DarK              & Dark subtracted \\
{\tt\_ff} & Flat Field        & Divided by the flat field \\
{\tt\_md} & Masked Deviants   & Deviant pixels from the neighbourhood (usually
                                3~$\sigma$ in 15$\times$15-pixel region) flagged as bad \\
{\tt\_nl} & Non-Linearity     & The standard non-linearity correction has
                                been applied \\
{\tt\_nm} & Normalised to Mode& Created in \htmlref{QUADRANT\_JITTER}{QUADRANT_JITTER}
                                to normalise frames combined to make the flat field \\
{\tt\_om} & Objects Masked    & This has sources masked with bad values
                                so that they do not bias the self flat field \\
{\tt\_qm} & Quadrant Masked   & One of the quadrants is masked with bad
                                pixels, created in QUADRANT\_JITTER \\
{\tt\_sbp}& Substitute Bad Pixels & Bad pixels replaced (needed for \PISA) \\
{\tt\_sc} & SCaled            & Data scaled to lie within the range of values
                                allowed by \xref{\PISA}{sun109}{} for the object 
                                masking \\
{\tt\_ss} & Sky Subtracted    & Global or local sky subtraction applied \\
{\tt\_trn}& TRaNsform         & The transformed or resampled data immediately 
                                prior to making a mosaic \\
\end{tabular}

\subsection*{UFTI Group suffices}%
\index{UFTI Group suffices}
\begin{tabular}{llp{80mm}}
\hline
  Suffix   & Stands for             & Description \\ \hline
{\tt\_mos} & Mosaic                 & Final mosaic \\
{\tt\_mu}  & Mosaic (Unfiltered)~~~ & Intermediate mosaic (could contain
                                       bad/hot pixels) \\
\end{tabular}

\newpage
\section{\xlabel{recipes}Recipes\label{recipes}}

The original set of recipes and names were prescribed in
G.S. Wright \& S.K. Leggett, 1997, {\em Scripts for UFTI}, 
\htmladdnormallink{orac009-ufts, v01}
{http://www.jach.hawaii.edu/JACpublic/UKIRT/software/orac/docs/orac009-ufts-1.html}.

\subsection{\xlabel{classified_recipes}Classified Recipes}

In hindsight you may decide that there was a better recipe for your
data than stored in the RECIPE header.  Here is a classified list so
that you can select an alternative.  Magnitudes and dimensions
apply to UFTI, except for the NOD recipes, whose magnitude ranges
are for IRCAM/TUFTI.  The magnitude ranges are courtesy of Sandy
Leggett.
\bigskip\bigskip

\begin{center}
\begin{tabular}{|l|l|p{75mm}|}
\multicolumn{3}{c}{\large{\bf Calibration Recipes}} \vspace*{1ex} \\
\hline
Recipe Name & Type of Data & Function and Comments \\ \hline
\htmlref{ARRAY\_TESTS}{ARRAY\_TESTS} & Array check &
   Calculates read noise and dark current \\ \hline
\htmlref{REDUCE\_DARK}{REDUCE\_DARK} & Dark &
   Files a frame as the current dark. \\ \hline
\htmlref{SKY\_FLAT}{SKY\_FLAT} & Flat &
   Creates and files a flat field derived from 5 jittered frames.
   Mostly for use with BRIGHT\_POINT\_SOURCE recipes. \\ \hline
\htmlref{SKY\_FLAT\_MASKED}{SKY\_FLAT\_MASKED} & Flat &
    As SKY\_FLAT but masks objects to give a better flat field. \\ \hline
\end{tabular}
\end{center}
\bigskip\bigskip

\begin{center}
\begin{tabular}{|l|p{24mm}|p{59mm}|}
\multicolumn{3}{c}{\large{\bf Very-bright-point-source recipes}} \vspace*{1ex} \\
\hline
Recipe Name & Type of Data & Function and Comments \\ \hline
\htmlref{BRIGHT\_POINT\_SOURCE}{BRIGHT\_POINT\_SOURCE} & 
   \mbox{$I$,$Z<13$}, \mbox{$JHK<9$}, and bright \mbox{$13<I$,$Z<17$}, \mbox{$9<JHK<15$} & 
   5-point jitter but would be useable as 3-point.  Requires a
   separate flat as the background is too low to self flat. \\ \hline
\htmlref{BRIGHT\_POINT\_SOURCE\_APHOT}{BRIGHT\_POINT\_SOURCE\_APHOT} & &
   As BRIGHT\_POINT\_SOURCE, but also performs aperture photometry of the
   source. \\ \hline
\end{tabular}
\end{center}

\newpage
\begin{center}
\begin{tabular}{|l|p{25mm}|p{59mm}|}
\multicolumn{3}{c}{\large{\bf Bright-point-source recipes}} \vspace*{1ex} \\
\hline
Recipe Name & Type of Data & Function and Comments \\ \hline
\htmlref{JITTER3\_SELF\_FLAT}{JITTER3\_SELF\_FLAT} &
   \mbox{$13<I$,$Z<17$}, \mbox{$9<JHK<15$} & 3-point standard jitter, self flats.
   Not recommended. \\ \hline
\htmlref{JITTER5\_SELF\_FLAT}{JITTER5\_SELF\_FLAT} & &
   5-point standard jitter, self flats. \\ \hline
\htmlref{JITTER5\_SELF\_FLAT\_APHOT}{JITTER5\_SELF\_FLAT\_APHOT} & &
   As JITTER5\_SELF\_FLAT, but also performs aperture photometry of the
   source. \\ \hline
\htmlref{JITTER5\_SELF\_FLAT\_NO\_MASK}{JITTER5\_SELF\_FLAT\_NO\_MASK} & &
   As JITTER5\_SELF\_FLAT but faster as it lacks object masking.  It
   only suitable for uncrowded fields. \\ \hline
\htmlref{SKY\_AND\_JITTER5}{SKY\_AND\_JITTER5} & &
   A sky frame and 5-point jitter.  The sky is subtracted from the
   target frame before flat fielding.  Requires a
   separate flat as the background is too low to self flat.
   No longer recommended as sky varies too quickly. \\ \hline
\htmlref{SKY\_AND\_JITTER5\_APHOT}{SKY\_AND\_JITTER5\_APHOT} & &
   As SKY\_AND\_JITTER5, but also performs aperture photometry of the
   source.  \\ \hline
\end{tabular}
\end{center}
\bigskip\bigskip

\begin{center}
\begin{tabular}{|l|p{26mm}|p{61mm}|}
\multicolumn{3}{c}{\large{\bf Faint-point-source recipes}} \vspace*{1ex} \\
\hline
Recipe Name & Type of Data & Function and Comments \\ \hline
\htmlref{JITTER9\_SELF\_FLAT}{JITTER9\_SELF\_FLAT} & Faint
   $I$,$Z>17$, $JHK>15$ & 9-point standard jitter, self flats. Otherwise
   known as JITTER\_SELF\_FLAT. \\ \hline
\htmlref{JITTER9\_SELF\_FLAT\_BASIC}{JITTER9\_SELF\_FLAT\_BASIC} & &
   Fastest JITTER9\_SELF\_FLAT recipe as it lacks object masking, automatic
   registration and resampling. \\ \hline
\htmlref{JITTER9\_SELF\_FLAT\_NO\_MASK}{JITTER9\_SELF\_FLAT\_NO\_MASK} & &
   As JITTER9\_SELF\_FLAT, but faster as it lacks object masking.  It
   only suitable for uncrowded fields. \\ \hline
\end{tabular}
\end{center}

\newpage
\begin{center}
\begin{tabular}{|l|p{24mm}|p{61mm}|}
\multicolumn{3}{c}{\large{\bf Point-source recipes---thermal}} \vspace*{1ex} \\
\hline
Recipe Name & Type of Data & Function and Comments \\ \hline
\htmlref{NOD4\_SELF\_FLAT\_NO\_MASK}{NOD4\_SELF\_FLAT\_NO\_MASK} & Bright $L<10$ &
   Nod 4-point jitter, self flats of differenced pairs of frames.  Has superior
   and fast sky subtraction.  No object masking. \\ \hline
\htmlref{NOD4\_SELF\_FLAT\_NO\_MASK\_APHOT}{NOD4\_SELF\_FLAT\_NO\_MASK\_APHOT} & &
   As NOD4\_SELF\_FLAT\_NO\_MASK, but also performs aperture photometry of the
   positive and negative sources.  \\ \hline
\htmlref{NOD8\_SELF\_FLAT\_NO\_MASK}{NOD8\_SELF\_FLAT\_NO\_MASK} & Faint $L>10$, all $M$ &
   Nod 8-point jitter, self flats of differenced pairs of frames.  Has superior
   sky subtraction.  No object masking. \\ \hline
\htmlref{NOD8\_SELF\_FLAT\_NO\_MASK\_APHOT}{NOD8\_SELF\_FLAT\_NO\_MASK\_APHOT} & &
   As NOD8\_SELF\_FLAT\_NO\_MASK, but also performs aperture photometry of the
   positive and negative sources.  \\ \hline
\end{tabular}
\end{center}

\begin{center}
\begin{tabular}{|l|p{32mm}|p{54mm}|}
\multicolumn{3}{c}{\large{\bf Extended-source recipes}} \vspace*{1ex} \\
\hline
Recipe Name & Type of Data & Function and Comments \\ \hline
\htmlref{QUADRANT\_JITTER}{QUADRANT\_JITTER} & Galaxies, quasars and
   nebulae of small ($<$45 arcsec) angular extent & 4-point jitter; masks
   the quadrant containing the target to make the flat, then
   masks the objects. \\ \hline
\htmlref{QUADRANT\_JITTER\_NO\_MASK}{QUADRANT\_JITTER\_NO\_MASK} & &
   As QUADRANT\_JITTER but without object masking. \\ \hline
\htmlref{QUADRANT\_JITTER\_BASIC}{QUADRANT\_JITTER\_BASIC} & &
   Fastest QUADRANT\_JITTER recipe as it lacks object masking, automatic
   registration and resampling. \\ \hline
\htmlref{EXTENDED\_3x3}{EXTENDED\_3x3} & Galaxies and nebulae with
   angular extent $<$2 arcminutes & Sky-subtracted 3$\times$3 grid
   mosaic.  Frames alternate between sky and target. \\ \hline
\htmlref{EXTENDED\_3x3\_BASIC}{EXTENDED\_3x3\_BASIC} & &
   As EXTENDED\_3x3 but lacks resampling and registers using telescope
   offsets. \\ \hline
\htmlref{EXTENDED\_5x5}{EXTENDED\_5x5} & Galaxies and nebulae with
   angular extent $<$3 arcminutes & Sky-subtracted 5$\times$5 grid 
   mosaic.  Frames alternate between sky and target. \\ \hline
\htmlref{EXTENDED\_5x5\_BASIC}{EXTENDED\_5x5\_BASIC} & &
   As EXTENDED\_5x5 but lacks resampling and registers using telescope
   offsets. \\ \hline
\end{tabular}
\end{center}

\begin{center}
\begin{tabular}{|l|p{14mm}|p{61mm}|}
\multicolumn{3}{c}{\large{\bf Moving (non-sidereal) source recipes}} \vspace*{1ex} \\
\hline
Recipe Name & Type of Data & Function and Comments \\ \hline
\htmlref{MOVING\_JITTER9\_SELF\_FLAT}{MOVING\_JITTER9\_SELF\_FLAT} &
   Minor planets, comets &
   As JITTER9\_SELF\_FLAT, but uses ephemeris data to track the non-sidereal source. \\ \hline
\htmlref{MOVING\_JITTER9\_SELF\_FLAT\_BASIC}{MOVING\_JITTER9\_SELF\_FLAT\_BASIC} &  &
   As JITTER9\_SELF\_FLAT\_BASIC, but uses ephemeris data to track the non-sidereal source.\\ \hline
\htmlref{JITTER9\_SELF\_FLAT\_TELE}{JITTER9\_SELF\_FLAT\_TELE} & &
   9-point standard jitter, using telescope offsets.  This is needed
   when the telescope has tracked on the non-sidereal target. \\ \hline
\end{tabular}
\end{center}

\subsection{Reference documentation}

The reference documentation which follows describes the UFTI recipes.
Since the IRCAM recipes are virtually identical, to list these, when
the UFTI recipes already have considerable repetition, was
profiligate. The differences to note are the frames have the {\tt ro}
prefix instead of {\tt f}, the names of group frames, i.e. mosaics,
begin with {\tt rg} instead of {\tt g} for UFTI.  The fields of view
as given in recipes like EXTENDED\_5x5 will be correspondingly
smaller, depending on the plate scale.  For IRCAM3 data without the
magnifier, the field of view is 72 arcseconds, compared with 91
arcseconds for UFTI's full-size frame.  IRCAM data will usually have a
non-linearity correction applied.

In the {\bf Configurable Steering Parameters} sections the defaults
appear at the end of the parameter's description between {\tt [~]}.

\newpage
\sstroutine{
   ARRAY\_TESTS
}{
   Calculates the ND\_STARE readout noise and dark current
}{
   \sstdescription{
      This script calculates the STARE and ND\_STARE readout noises,
      and the dark current from a series of engineering data taken
      with the EXEC called ARRAY\_TESTS.  The results are compared with
      the nominal values, and you are notified whether or not the values
      obtained are within limits.  At UKIRT, the results are also logged
      to an engineering file for archival purposes.
   }
   \sstnotes{
      \sstitemlist{

         \sstitem
         Intermediate frames are deleted.

         \sstitem
         The engineering log contains the UT date and time, the ND\_STARE
         readout noise, and the dark current.  The results are normally
         appended to the log.  If for some reason it does not exist, a new
         log is created containing the column headings.

         \sstitem
         Multiple array tests are permitted.  A new set of results is
         reported and logged for each cycle.
      }
   }
   \sstdiytopic{
      Output Data
   }{
      \sstitemlist{

         \sstitem
         The engineering log: {\tt /ukirt\_sw/logs/ufti\_array\_tests.log}.
      }
   }
   \sstimplementationstatus{
      \sstitemlist{

         \sstitem
         The processing engines are from the Starlink package \xref{\KAPPA}{sun95}{}.

         \sstitem
         Uses the Starlink NDF format.
      }
   }
}
\newpage
\sstroutine{
   BRIGHT\_POINT\_SOURCE
}{
   Reduces a bright-point-source photometry observation
}{
   \sstdescription{
      This recipe reduces a ``bright standard'' photometry observation
      with UFTI data.  It takes an UFTI observation comprising a
      series of object frames and a dark frame, and a predetermined
      flat-field frame to make a calibrated, trimmed mosaic
      automatically.

      This recipe performs bad-pixel masking, null debiassing, dark
      subtraction, flat-field division, feature detection and matching
      between object frames, and resampling.  See the notes for details.

      As the name implies, it is intended for bright point sources,
      such as standard stars.
   }
   \sstnotes{
      \sstitemlist{

         \sstitem
         You may use SKY\_FLAT or SKY\_FLAT\_MASKED to make the flat field.

         \sstitem
         The bad pixel mask applied is {\tt\$ORAC\_DATA\_CAL/bpm}.

         \sstitem
         Where automatic registration is not possible, the recipe matches
         the centroid of central source, and should that fail, it resorts
         to using the telescope offsets transformed to pixels.

         \sstitem
         The resampling applies non-integer shifts of origin using
         bilinear interpolation.  There is no rotation to align the
         Cartesian axes with the cardinal directions.

         \sstitem
         The recipe makes the mosaic by applying offsets in intensity to
         give the most consistent result amongst the overlapping regions.
         The mosaic is trimmed to the dimensions of an input frame.  The
         mosaic is not normalised by its exposure time (that being the
         exposure time of a single frame).

         \sstitem
         For each cycle of five, the recipe creates a mosaic, which
         is then added into a master mosaic of improving signal to noise.
         The exposure time is also summed and stored in the mosaic's
         EXP\_TIME header.  Likewise the end airmass header, AMEND, is
         updated to match that of the last-observed frame contributing to
         the mosaic.

         \sstitem
         Intermediate frames are deleted except for the flat-fielded
         ({\tt\_ff} suffix) frames.

         \sstitem
         Sub-arrays are supported.
      }
   }
   \sstdiytopic{
      Output Data
   }{
      \sstitemlist{

         \sstitem
         The resultant mosaic in {\tt g$<$date$>$\_$<$group\_number$>$\_mos}.

         \sstitem
         The individual flat-fielded frames in {\tt f$<$date$>$\_$<$obs\_number$>$\_ff}.
      }
   }
   \sstparameters{
      \sstsubsection{
         NUMBER = INTEGER
      }{
         The number of frames in the jitter pattern. {\tt[5]}
      }
   }
   \sstdiytopic{
      Timing
   }{
      This is a fast recipe and will keep pace with observations.
   }
   \sstdiytopic{
      Related Recipes
   }{
      \htmlref{BRIGHT\_POINT\_SOURCE\_APHOT}{BRIGHT\_POINT\_SOURCE\_APHOT},
      \htmlref{JITTER5\_SELF\_FLAT\_APHOT}{JITTER5\_SELF\_FLAT\_APHOT},
      \htmlref{SKY\_FLAT}{SKY\_FLAT},\\
      \htmlref{SKY\_FLAT\_MASKED}{SKY\_FLAT\_MASKED}.
   }
   \sstimplementationstatus{
      \sstitemlist{

         \sstitem
         The processing engines are from the Starlink packages: \xref{\CCDPACK}{sun139}{}
         and \xref{\KAPPA}{sun95}{}.

         \sstitem
         Uses the Starlink NDF format.

         \sstitem
         History is recorded within the data files.

         \sstitem
         The title of the data is propagated through intermediate files
         to the mosaic.

         \sstitem
         Error propagation is not used.
      }
   }
}
\newpage
\sstroutine{
   BRIGHT\_POINT\_SOURCE\_APHOT
}{
   Reduces a bright-point-source photometry observation and performs
   aperture photometry
}{
   \sstdescription{
      This recipe reduces a ``bright standard'' photometry observation
      with UFTI data.  It takes an UFTI observation comprising a
      series of object frames and a dark frame with a predetermined
      flat-field frame to make a calibrated, trimmed mosaic
      automatically.

      This recipe performs a null debiassing, bad-pixel masking, dark
      subtraction, flat-field division, feature detection and matching
      between object frames, and resampling.  See the notes for details.

      Photometry of the point source using a fixed 6-arcsecond aperture
      is calculated for each jitter frame and the mosaic.  The results
      appear in {\tt\$ORAC\_DATA\_OUT/aphot\_results.txt} in the form of 
      a Starlink \xref{small text list}{sun190}{STLREF}.  The analysis
      of each star is appended to this file.

      As the name implies, it is intended for bright point sources,
      such as standard stars.
   }
   \sstnotes{
      \sstitemlist{

         \sstitem
         You may use \htmlref{SKY\_FLAT}{SKY\_FLAT} or \htmlref{SKY\_FLAT\_MASKED}{SKY\_FLAT\_MASKED} to make the flat field.

         \sstitem
         The bad pixel mask applied is {\tt\$ORAC\_DATA\_CAL/bpm}.

         \sstitem
         Where automatic registration is not possible, the recipe matches
         the centroid of central source, and should that fail, it resorts
         to using the telescope offsets transformed to pixels.

         \sstitem
         The resampling applies non-integer shifts of origin using
         bilinear interpolation.  There is no rotation to align the
         Cartesian axes with the cardinal directions.

         \sstitem
         The recipe makes the mosaic by applying offsets in intensity to
         give the most consistent result amongst the overlapping regions.
         The mosaic is trimmed to the dimensions of an input frame.  The
         mosaic is not normalised by its exposure time (that being the
         exposure time of a single frame).

         \sstitem
         For each cycle of five, the recipe creates a mosaic, which
         is then added into a master mosaic of improving signal to noise.
         The exposure time is also summed and stored in the mosaic{\tt '}s
         EXP\_TIME header.  Likewise the end airmass header, AMEND, is
         updated to match that of the last-observed frame contributing to
         the mosaic.

         \sstitem
         The photometry tabulation includes the file name, source name,
         time, filter, airmass, the catalogue magnitude and estimates of
         the zero-point with and without the application of a mean
         extinction.  There are headings at the top of each column.

         \sstitem
         Intermediate frames are deleted except for the flat-fielded ({\tt\_ff}
         suffix) frames.

         \sstitem
         Sub-arrays are supported.
      }
   }
   \sstdiytopic{
      Output Data
   }{
      \sstitemlist{

         \sstitem
         The resultant mosaic in {\tt g$<$date$>$\_$<$group\_number$>$\_mos}.

         \sstitem
         The individual flat-fielded frames in {\tt f$<$date$>$\_$<$obs\_number$>$\_ff}.

         \sstitem
         Results tabulation to log {\tt \$ORAC\_DATA\_OUT/aphot\_results.txt}.
      }
   }
   \sstparameters{
      \sstsubsection{
         NUMBER = INTEGER
      }{
         The number of frames in the jitter pattern. {\tt[5]}
      }
   }
   \sstdiytopic{
      Timing
   }{
      This is a fast recipe and will keep pace with observations.
   }
   \sstdiytopic{
      Related Recipes
   }{
      \htmlref{BRIGHT\_POINT\_SOURCE}{BRIGHT\_POINT\_SOURCE},
      \htmlref{SKY\_FLAT}{SKY\_FLAT},
      \htmlref{SKY\_FLAT\_MASKED}{SKY\_FLAT\_MASKED}.
   }
   \sstimplementationstatus{
      \sstitemlist{

         \sstitem
         The processing engines are from the Starlink packages: \xref{\CCDPACK}{sun139}{},
         \xref{\KAPPA}{sun95}{}, and \xref{\PHOTOM}{sun45}{}.

         \sstitem
         Uses the Starlink NDF format.

         \sstitem
         History is recorded within the data files.

         \sstitem
         The title of the data is propagated through intermediate files
         to the mosaic.

         \sstitem
         Error propagation is not used.
      }
   }
}
\newpage
\sstroutine{
   EXTENDED\_3x3
}{
   Extended source standard reduction using interpolated sky subtraction
}{
   \sstdescription{
      This recipe reduces an extended source using UFTI data.  The data
      comprise alternating blank-sky and target frames commencing and
      ending with a blank sky.  The target frames are arranged in an
      overlapping (30-50\%) grid of 3$\times$3 frames from which the recipe
      makes a sky-subtracted untrimmed mosaic automatically.

      The script performs bad-pixel masking, null debiassing, dark
      subtraction, flat-field division, sky subtraction, registration,
      resampling, and mosaicking.  The notes give more details.

      It is suitable for extended objects up to 2 arcminutes across.
   }
   \sstnotes{
      \sstitemlist{

         \sstitem
         The bad pixel mask applied is {\tt\$ORAC\_DATA\_CAL/bpm}.

         \sstitem
         The flat field is derived from the sky frames as follows.  The
         mode (sigma-clipped mean) is used to offset each sky frame's mode
         to that of the first sky frame.  The corrected sky frames are
         combined pixel by pixel using a median of the values in each
         frame.  The resultant frame is normalised by its median to form
         the flat field.  This frame median is subtracted from the source
         frames after they have been flat-fielded.  A flat field is created
         for each row of the grid of target frames, and applied only to
         that row of target frames.

         \sstitem
         The sky subtraction comes from linear interpolation of the sky
         modal values of the two sky frames which immediately bracket the
         target frame.

         \sstitem
         Registration is performed using common point sources in the
         overlap regions.  If the recipe cannot identify sufficient common
         objects, it then tries the crosshead offsets.  If these are null,
         the script resorts to the telescope offsets.

         \sstitem
         The resampling applies non-integer shifts of origin using
         bilinear interpolation.  There is no rotation to align the
         Cartesian axes with the cardinal directions.

         \sstitem
         The recipe makes the mosaics by applying offsets in intensity
         to give the most consistent result amongst the overlapping regions.
         The noise will be greater in the mosaic's peripheral areas, having
         received less exposure time.  The mosaic is not normalised by its
         exposure time (that being the exposure time of a single frame).

         \sstitem
         Mosaics are made and displayed for each row, except the last.
         At the end of each cycle of 19 frames the full mosaic of nine target
         frames is created and displayed instead.  On the second and
         subsequent cycles the full mosaic is added into a master mosaic of
         improving signal to noise.  The exposure time is also summed and
         stored in the mosaic's EXP\_TIME header.  Likewise the end airmass
         header, AMEND, is updated to match that of the last-observed frame
         contributing to the mosaic.

         \sstitem
         Intermediate frames are deleted except for the flat-fielded ({\tt\_ff}
         suffix) frames.
      }
   }
   \sstimplementationstatus{
      \sstitemlist{

         \sstitem
         The processing engines are from the Starlink packages: \xref{\CCDPACK}{sun139}{},
         \xref{\KAPPA}{sun95}{}, and  \xref{\PISA}{sun109}{}.

         \sstitem
         Uses the Starlink NDF format.

         \sstitem
         History is recorded within the data files.

         \sstitem
         The title of the data is propagated through intermediate files
         to the mosaic.

         \sstitem
         Error propagation is not used.
      }
   }
}
\newpage
\sstroutine{
   EXTENDED\_3x3\_BASIC
}{
   Extended source standard reduction using interpolated sky subtraction
}{
   \sstdescription{
      This recipe reduces an extended source using UFTI data.  The data
      comprise alternating blank-sky and target frames commencing and
      ending with a blank sky.  The target frames are arranged in an
      overlapping (30-50\%) grid of 3$\times$3 frames from which the recipe
      makes a sky-subtracted untrimmed mosaic automatically.

      The script performs bad-pixel masking, null debiassing, dark
      subtraction, flat-field division, sky subtraction, registration
      using telescope offsets, and mosaicking.  The notes give more
      details.

      It is suitable for extended objects up to 2 arcminutes across.
   }
   \sstnotes{
      \sstitemlist{

         \sstitem
         The bad pixel mask applied is {\tt\$ORAC\_DATA\_CAL/bpm}.

         \sstitem
         The flat field is derived from the sky frames as follows.  The
         mode (sigma-clipped mean) is used to offset each sky frame's mode
         to that of the first sky frame.  The corrected sky frames are
         combined pixel by pixel using a median of the values in each
         frame.  The resultant frame is normalised by its median to form
         the flat field.  This frame median is subtracted from the source
         frames after they have been flat-fielded.  A flat field is created
         for each row of the grid of target frames, and applied only to
         that row of target frames.

         \sstitem
         The sky subtraction comes from linear interpolation of the sky
         modal values of the two sky frames which immediately bracket the
         target frame.

         \sstitem
         Registration is performed using the telescope offsets
         transformed to pixels.

         \sstitem
         There is no resampling, merely integer shifts of origin.

         \sstitem
         The recipe makes the mosaics by applying offsets in intensity
         to give the most consistent result amongst the overlapping regions.
         The noise will be greater in the mosaic's peripheral areas, having
         received less exposure time.  The mosaic is not normalised by its
         exposure time (that being the exposure time of a single frame).

         \sstitem
         Mosaics are made and displayed for each row, except the last.
         At the end of each cycle of 19 frames the full mosaic of nine target
         frames is created and displayed instead.  On the second and
         subsequent cycles the full mosaic is added into a master mosaic of
         improving signal to noise.  The exposure time is also summed and
         stored in the mosaic's EXP\_TIME header.  Likewise the end airmass
         header, AMEND, is updated to match that of the last-observed frame
         contributing to the mosaic.

         \sstitem
         Intermediate frames are deleted except for the flat-fielded ({\tt\_ff}
         suffix) frames.
      }
   }
   \sstdiytopic{
      Output Data
   }{
      \sstitemlist{

         \sstitem
         The full mosaic in {\tt g$<$date$>$\_$<$group\_number$>$\_mos}.

         \sstitem
         A mosaic for each row in {\tt g$<$date$>$\_$<$group\_number$>$\_mos$<$row\_number$>$},
         where {\tt$<$row\_number$>$} is 0 or 1.

         \sstitem
         The individual flat-fielded frames in {\tt f$<$date$>$\_$<$obs\_number$>$\_ff}.
      }
   }
   \sstparameters{
      \sstsubsection{
         NROW = INTEGER
      }{
         The number of target frames in a row of the mosaic.  Its
         minimum is 3 because this number of blank skies are needed to
         form a flat field. {\tt[3]}
      }
      \sstsubsection{
         NCOL = INTEGER
      }{
         The number of target frames in a column of the mosaic.  Its
         minimum is 2. {\tt[3]}
      }
   }
   \sstdiytopic{
      Related Recipes
   }{
      \htmlref{EXTENDED\_3x3}{EXTENDED\_3x3},
      \htmlref{EXTENDED\_5x5\_BASIC}{EXTENDED\_5x5\_BASIC},
      \htmlref{QUADRANT\_JITTER}{QUADRANT\_JITTER}.
   }
   \sstimplementationstatus{
      \sstitemlist{

         \sstitem
         The processing engines are from the Starlink packages: \xref{\CCDPACK}{sun139}{},
         \xref{\KAPPA}{sun95}{}, and  \xref{\PISA}{sun109}{}.

         \sstitem
         Uses the Starlink NDF format.

         \sstitem
         History is recorded within the data files.

         \sstitem
         The title of the data is propagated through intermediate files
         to the mosaic.

         \sstitem
         Error propagation is not used.
      }
   }
}
\newpage
\sstroutine{
   EXTENDED\_5x5
}{
   Extended source standard reduction using interpolated sky subtraction
}{
   \sstdescription{
      This recipe reduces an extended source using UFTI data.  The data
      comprise alternating blank-sky and target frames commencing and
      ending with a blank sky.  The target frames are arranged in an
      overlapping (30-50\%) grid of 5$\times$5 frames from which the recipe
      makes a sky-subtracted untrimmed mosaic automatically.

      The script performs bad-pixel masking, null debiassing, dark
      subtraction, flat-field division, sky subtraction, registration,
      resampling, and mosaicking.  The notes give more details.

      It is suitable for extended objects up to 3 arcminutes across.
   }
   \sstnotes{
      \sstitemlist{

         \sstitem
         The bad pixel mask applied is {\tt\$ORAC\_DATA\_CAL/bpm}.

         \sstitem
         The flat field is derived from the sky frames as follows.  The
         mode (sigma-clipped mean) is used to offset each sky frame's mode
         to that of the first sky frame.  The corrected sky frames are
         combined pixel by pixel using a median of the values in each
         frame.  The resultant frame is normalised by its median to form
         the flat field.  This frame median is subtracted from the source
         frames after they have been flat-fielded.  A flat field is created
         for each row of the grid of target frames, and applied only to
         that row of target frames.

         \sstitem
         The sky subtraction comes from linear interpolation of the sky
         modal values of the two sky frames which immediately bracket the
         target frame.

         \sstitem
         Registration is performed using common point sources in the
         overlap regions.  If the recipe cannot identify sufficient common
         objects, it then tries the crosshead offsets.  If these are null,
         the script resorts to the telescope offsets.

         \sstitem
         The resampling applies non-integer shifts of origin using
         bilinear interpolation.  There is no rotation to align the
         Cartesian axes with the cardinal directions.

         \sstitem
         The recipe makes the mosaics by applying offsets in intensity
         to give the most consistent result amongst the overlapping regions.
         The noise will be greater in the mosaic's peripheral areas, having
         received less exposure time.  The mosaic is not normalised by its
         exposure time (that being the exposure time of a single frame).

         \sstitem
         Mosaics are made and displayed for each row, except the last.
         At the end of each cycle of 51 frames the full mosaic of 25 target
         frames is created and displayed instead.  On the second and
         subsequent cycles the full mosaic is added into a master mosaic of
         improving signal to noise.  The exposure time is also summed and
         stored in the mosaic's EXP\_TIME header.  Likewise the end airmass
         header, AMEND, is updated to match that of the last-observed frame
         contributing to the mosaic.

         \sstitem
         Intermediate frames are deleted except for the flat-fielded ({\tt\_ff}
         suffix) frames.
      }
   }
   \sstimplementationstatus{
      \sstitemlist{

         \sstitem
         The processing engines are from the Starlink packages: \xref{\CCDPACK}{sun139}{},
         \xref{\KAPPA}{sun95}{}, and  \xref{\PISA}{sun109}{}.

         \sstitem
         Uses the Starlink NDF format.

         \sstitem
         History is recorded within the data files.

         \sstitem
         The title of the data is propagated through intermediate files
         to the mosaic.

         \sstitem
         Error propagation is not used.
      }
   }
}
\newpage
\sstroutine{
   EXTENDED\_5x5\_BASIC
}{
   Extended source standard reduction using interpolated sky subtraction
}{
   \sstdescription{
      This recipe reduces an extended source using UFTI data.  The data
      comprise alternating blank-sky and target frames commencing and
      ending with a blank sky.  The target frames are arranged in an
      overlapping (30-50\%) grid of 5$\times$5 frames from which the recipe
      makes a sky-subtracted untrimmed mosaic automatically.

      The script performs bad-pixel masking, null debiassing, dark
      subtraction, flat-field division, sky subtraction, registration
      using telescope offsets, and mosaicking.  The notes give more
      details.

      It is suitable for extended objects up to 3 arcminutes across.
   }
   \sstnotes{
      \sstitemlist{

         \sstitem
         The bad pixel mask applied is {\tt\$ORAC\_DATA\_CAL/bpm}.

         \sstitem
         The flat field is derived from the sky frames as follows.  The
         mode (sigma-clipped mean) is used to offset each sky frame's mode
         to that of the first sky frame.  The corrected sky frames are
         combined pixel by pixel using a median of the values in each
         frame.  The resultant frame is normalised by its median to form
         the flat field.  This frame median is subtracted from the source
         frames after they have been flat-fielded.  A flat field is created
         for each row of the grid of target frames, and applied only to
         that row of target frames.

         \sstitem
         The sky subtraction comes from linear interpolation of the sky
         modal values of the two sky frames which immediately bracket the
         target frame.

         \sstitem
         Registration is performed using the telescope offsets
         transformed to pixels.

         \sstitem
         There is no resampling, merely integer shifts of origin.

         \sstitem
         The recipe makes the mosaics by applying offsets in intensity
         to give the most consistent result amongst the overlapping regions.
         The noise will be greater in the mosaic's peripheral areas, having
         received less exposure time.  The mosaic is not normalised by its
         exposure time (that being the exposure time of a single frame).

         \sstitem
         Mosaics are made and displayed for each row, except the last.
         At the end of each cycle of 51 frames the full mosaic of 25 target
         frames is created and displayed instead.  On the second and
         subsequent cycles the full mosaic is added into a master mosaic of
         improving signal to noise.  The exposure time is also summed and
         stored in the mosaic's EXP\_TIME header.  Likewise the end airmass
         header, AMEND, is updated to match that of the last-observed frame
         contributing to the mosaic.

         \sstitem
         Intermediate frames are deleted except for the flat-fielded ({\tt\_ff}
         suffix) frames.
      }
   }
   \sstdiytopic{
      Output Data
   }{
      \sstitemlist{

         \sstitem
         The full mosaic in {\tt g$<$date$>$\_$<$group\_number$>$\_mos}.

         \sstitem
         A mosaic for each row in {\tt g$<$date$>$\_$<$group\_number$>$\_mos$<$row\_number$>$},
         where {\tt$<$row\_number$>$} is 0 to 3.

         \sstitem
         The individual flat-fielded frames in {\tt f$<$date$>$\_$<$obs\_number$>$\_ff}.
      }
   }
   \sstparameters{
      \sstsubsection{
         NROW = INTEGER
      }{
         The number of target frames in a row of the mosaic.  Its
         minimum is 3 because this number of blank skies are needed to
         form a flat field. {\tt[5]}
      }
      \sstsubsection{
         NCOL = INTEGER
      }{
         The number of target frames in a column of the mosaic.  Its
         minimum is 2. {\tt[5]}
      }
   }
   \sstdiytopic{
      Related Recipes
   }{
      \htmlref{EXTENDED\_3x3}{EXTENDED\_3x3},
      \htmlref{EXTENDED\_3x3\_BASIC}{EXTENDED\_3x3\_BASIC},
      \htmlref{EXTENDED\_5x5}{EXTENDED\_5x5},
      \htmlref{QUADRANT\_JITTER}{QUADRANT\_JITTER}.
   }
   \sstimplementationstatus{
      \sstitemlist{

         \sstitem
         The processing engines are from the Starlink packages: \xref{\CCDPACK}{sun139}{},
         \xref{\KAPPA}{sun95}{}, and  \xref{\PISA}{sun109}{}.

         \sstitem
         Uses the Starlink NDF format.

         \sstitem
         History is recorded within the data files.

         \sstitem
         The title of the data is propagated through intermediate files
         to the mosaic.

         \sstitem
         Error propagation is not used.
      }
   }
}
\newpage
\sstroutine{
   JITTER3\_SELF\_FLAT
}{
   Reduces a 3-point ``standard jitter'' photometry observation using
   object masking
}{
   \sstdescription{
      This script reduces a ``standard jitter'' photometry observation
      with UFTI data.  It takes an UFTI observation comprising three
      object frames and a dark frame to make a calibrated, untrimmed
      mosaic automatically.

      It performs a null debiassing, bad-pixel masking, dark
      subtraction, flat-field creation and division, feature
      detection and matching between object frames, and resampling.

      This recipe works well for faint sources and for moderately
      crowded fields.
   }
   \sstnotes{
      \sstitemlist{

         \sstitem
         The bad pixel mask applied is {\tt\$ORAC\_DATA\_CAL/bpm}.

         \sstitem
         The flat field is created iteratively.  First an approximate
         flat-field is created by combining normalised object frames using
         the median at each pixel.  This flat field is applied to the object
         frames.  Sources within the flat-fielded frames are detected, and
         masked in the dark-subtracted frames.  The first stage is repeated
         but applied to the masked frames to create the final flat field.

         \sstitem
         Registration is performed using common point sources in the
         overlap regions.  If the recipe cannot identify sufficient common
         objects, the script resorts to using the telescope offsets
         transformed to pixels.

         \sstitem
         The resampling applies non-integer shifts of origin using
         bilinear interpolation.  There is no rotation to align the
         Cartesian axes with the cardinal directions.

         \sstitem
         The recipe makes the mosaics by applying offsets in intensity
         to give the most consistent result amongst the overlapping regions.
         The mosaic is not trimmed to the dimensions of a single frame, thus
         the noise will be greater in the peripheral areas having received
         less exposure time.  The mosaic is not normalised by its exposure
         time (that being the exposure time of a single frame).

         \sstitem
         For each cycle of three, the recipe creates a mosaic, which has
         its bad pixels filled and is then added into a master mosaic of
         improving signal to noise.  The exposure time is also summed and
         stored in the mosaic's EXP\_TIME header.  Likewise the end airmass
         header, AMEND, is updated to match that of the last-observed frame
         contributing to the mosaic.

         \sstitem
         Intermediate frames are deleted except for the flat-fielded ({\tt\_ff}
         suffix) frames.
      }
   }
   \sstdiytopic{
      Output Data
   }{
      \sstitemlist{

         \sstitem
         The integrated mosaic in {\tt g$<$date$>$\_$<$group\_number$>$\_mos}.

         \sstitem
         A mosaic for each cycle of four in
         {\tt g$<$date$>$\_$<$group\_number$>$\_mos\_$<$cycle\_number$>$},
         where {\tt$<$cycle\_number$>$} counts from 0.

         \sstitem
         The individual flat-fielded frames in {\tt f$<$date$>$\_$<$obs\_number$>$\_ff}.

         \sstitem
         The created flat fields in {\tt flat\_$<$group\_number$>$subgrp} for the
         first or only cycle, and {\tt flat\_$<$group\_number$>$subgrp\_$<$cycle\_number$>$}
         for subsequent cycles.
      }
   }
   \sstparameters{
      \sstsubsection{
         NUMBER = INTEGER
      }{
         The number of frames in the jitter pattern. {\tt[3]}
      }
   }
   \sstdiytopic{
      Related Recipes
   }{
      \htmlref{JITTER5\_SELF\_FLAT}{JITTER5\_SELF\_FLAT},
      \htmlref{JITTER5\_SELF\_FLAT\_APHOT}{JITTER5\_SELF\_FLAT\_APHOT},\\
      \htmlref{JITTER5\_SELF\_FLAT\_NO\_MASK}{JITTER5\_SELF\_FLAT\_NO\_MASK},
      \htmlref{QUADRANT\_JITTER}{QUADRANT\_JITTER},\\
      \htmlref{QUADRANT\_JITTER\_NO\_MASK}{QUADRANT\_JITTER\_NO\_MASK}.
   }
   \sstimplementationstatus{
      \sstitemlist{

         \sstitem
         The processing engines are from the Starlink packages: \xref{\CCDPACK}{sun139}{},
         \xref{\KAPPA}{sun95}{}, and  \xref{\PISA}{sun109}{}.

         \sstitem
         Uses the Starlink NDF format.

         \sstitem
         History is recorded within the data files.

         \sstitem
         The title of the data is propagated through intermediate files
         to the mosaic.

         \sstitem
         Error propagation is not used.
      }
   }
}
\newpage
\sstroutine{
   JITTER5\_SELF\_FLAT
}{
   Reduces a 5-point ``standard jitter'' photometry observation using
   object masking
}{
   \sstdescription{
      This script reduces a ``standard jitter'' photometry observation
      with UFTI data.  It takes an UFTI observation comprising five
      object frames and a dark frame to make a calibrated, untrimmed
      mosaic automatically.

      It performs a null debiassing, bad-pixel masking, dark
      subtraction, flat-field creation and division, feature
      detection and matching between object frames, and resampling.

      This recipe works well for faint sources and for moderately
      crowded fields.
   }
   \sstnotes{
      \sstitemlist{

         \sstitem
         The bad pixel mask applied is {\tt\$ORAC\_DATA\_CAL/bpm}.

         \sstitem
         The flat field is created iteratively.  First an approximate
         flat-field is created by combining normalised object frames using
         the median at each pixel.  This flat field is applied to the object
         frames.  Sources within the flat-fielded frames are detected, and
         masked in the dark-subtracted frames.  The first stage is repeated
         but applied to the masked frames to create the final flat field.

         \sstitem
         Registration is performed using common point sources in the
         overlap regions.  If the recipe cannot identify sufficient common
         objects, the script resorts to using the telescope offsets
         transformed to pixels.

         \sstitem
         The resampling applies non-integer shifts of origin using
         bilinear interpolation.  There is no rotation to align the
         Cartesian axes with the cardinal directions.

         \sstitem
         The recipe makes the mosaics by applying offsets in intensity
         to give the most consistent result amongst the overlapping regions.
         The mosaic is not trimmed to the dimensions of a single frame, thus
         the noise will be greater in the peripheral areas having received
         less exposure time.

         \sstitem
         For each cycle of five, the recipe creates a mosaic, which has
         its bad pixels filled and is then added into a master mosaic of
         improving signal to noise.  The exposure time is also summed and
         stored in the mosaic's EXP\_TIME header.  Likewise the end airmass
         header, AMEND, is updated to match that of the last-observed frame
         contributing to the mosaic.

         \sstitem
         Intermediate frames are deleted except for the flat-fielded ({\tt\_ff}
         suffix) frames.
      }
   }
   \sstdiytopic{
      Output Data
   }{
      \sstitemlist{

         \sstitem
         The integrated mosaic in {\tt g$<$date$>$\_$<$group\_number$>$\_mos}.

         \sstitem
         A mosaic for each cycle of four in
         {\tt g$<$date$>$\_$<$group\_number$>$\_mos\_$<$cycle\_number$>$}, where {\tt$<$cycle\_number$>$}
         counts from 0.

         \sstitem
         The individual flat-fielded frames in {\tt f$<$date$>$\_$<$obs\_number$>$\_ff}.

         \sstitem
         The created flat fields in {\tt flat\_$<$group\_number$>$subgrp} for the
         first or only cycle, and {\tt flat\_$<$group\_number$>$subgrp\_$<$cycle\_number$>$}
         for subsequent cycles.
      }
   }
   \sstparameters{
      \sstsubsection{
         NUMBER = INTEGER
      }{
         The number of frames in the jitter pattern. {\tt[5]}
      }
   }
   \sstdiytopic{
      Related Recipes
   }{
      \htmlref{JITTER5\_SELF\_FLAT\_APHOT}{JITTER5\_SELF\_FLAT\_APHOT},
      \htmlref{JITTER5\_SELF\_FLAT\_NO\_MASK}{JITTER5\_SELF\_FLAT\_NO\_MASK},\\
      \htmlref{JITTER9\_SELF\_FLAT}{JITTER9\_SELF\_FLAT},
      \htmlref{JITTER9\_SELF\_FLAT\_NO\_MASK}{JITTER9\_SELF\_FLAT\_NO\_MASK},
      \htmlref{QUADRANT\_JITTER}{QUADRANT\_JITTER},\\
      \htmlref{QUADRANT\_JITTER\_NO\_MASK}{QUADRANT\_JITTER\_NO\_MASK}.
   }
   \sstimplementationstatus{
      \sstitemlist{

         \sstitem
         The processing engines are from the Starlink packages: \xref{\CCDPACK}{sun139}{},
         \xref{\KAPPA}{sun95}{}, and  \xref{\PISA}{sun109}{}.

         \sstitem
         Uses the Starlink NDF format.

         \sstitem
         History is recorded within the data files.

         \sstitem
         The title of the data is propagated through intermediate files
         to the mosaic.

         \sstitem
         Error propagation is not used.
      }
   }
}
\newpage
\sstroutine{
   JITTER5\_SELF\_FLAT\_APHOT
}{
   Reduces a 5-point ``standard jitter'' photometry observation using
   object masking, and performs aperture photometry
}{
   \sstdescription{
      This script reduces a ``standard jitter'' photometry observation
      with UFTI data.  It takes an UFTI observation comprising five
      object frames and a dark frame to make a calibrated, untrimmed
      mosaic automatically.

      It performs a null debiassing, bad-pixel masking, dark
      subtraction, flat-field creation and division, feature
      detection and matching between object frames, and resampling.

      Photometry of the point source using a fixed 6-arcsecond aperture
      is calculated for each jitter frame and the mosaic.  The results
      appear in {\tt\$ORAC\_DATA\_OUT/aphot\_results.txt} in the form of a Starlink
      \xref{small text list}{sun190}{STLREF}.  The analysis of each star is appended to this file.

      This recipe works well for faint sources in moderately crowded fields.
   }
   \sstnotes{
      \sstitemlist{

         \sstitem
         The bad pixel mask applied is {\tt\$ORAC\_DATA\_CAL/bpm}.

         \sstitem
         The flat field is created iteratively.  First an approximate
         flat-field is created by combining normalised object frames using
         the median at each pixel.  This flat field is applied to the object
         frames.  Sources within the flat-fielded frames are detected, and
         masked in the dark-subtracted frames.  The first stage is repeated
         but applied to the masked frames to create the final flat field.

         \sstitem
         Registration is performed using common point sources in the
         overlap regions.  If the recipe cannot identify sufficient common
         objects, the script resorts to using the telescope offsets
         transformed to pixels.

         \sstitem
         The resampling applies non-integer shifts of origin using
         bilinear interpolation.  There is no rotation to align the
         Cartesian axes with the cardinal directions.

         \sstitem
         The recipe makes the mosaics by applying offsets in intensity
         to give the most consistent result amongst the overlapping regions.
         The mosaic is not trimmed to the dimensions of a single frame, thus
         the noise will be greater in the peripheral areas having received
         less exposure time.  The mosaic is not normalised by its exposure
         time (the exposure time of a single frame).

         \sstitem
         For each cycle of five, the recipe creates a mosaic, which has
         its bad pixels filled and is then added into a master mosaic of
         improving signal to noise.  The exposure time is also summed and
         stored in the mosaic's EXP\_TIME header.  Likewise the end airmass
         header, AMEND, is updated to match that of the last-observed frame
         contributing to the mosaic.

         \sstitem
         Intermediate frames are deleted except for the flat-fielded ({\tt\_ff}
         suffix) frames.
      }
   }
   \sstdiytopic{
      Output Data
   }{
      \sstitemlist{

         \sstitem
         The integrated mosaic in {\tt g$<$date$>$\_$<$group\_number$>$\_mos}.

         \sstitem
         A mosaic for each cycle of four in
         {\tt g$<$date$>$\_$<$group\_number$>$\_mos\_$<$cycle\_number$>$}, where {\tt$<$cycle\_number$>$}
         counts from 0.

         \sstitem
         The individual flat-fielded frames in {\tt f$<$date$>$\_$<$obs\_number$>$\_ff}.

         \sstitem
         The created flat fields in {\tt flat\_$<$group\_number$>$subgrp} for the
         first or only cycle, and {\tt flat\_$<$group\_number$>$subgrp\_$<$cycle\_number$>$}
         for subsequent cycles.
      }
   }
   \sstparameters{
      \sstsubsection{
         NUMBER = INTEGER
      }{
         The number of frames in the jitter pattern. {\tt[5]}
      }
   }
   \sstdiytopic{
      Related Recipes
   }{
      \htmlref{JITTER5\_SELF\_FLAT}{JITTER5\_SELF\_FLAT},
      \htmlref{JITTER5\_SELF\_FLAT\_NO\_MASK}{JITTER5\_SELF\_FLAT\_NO\_MASK},
      \htmlref{JITTER9\_SELF\_FLAT}{JITTER9\_SELF\_FLAT}, \\
      \htmlref{JITTER9\_SELF\_FLAT\_NO\_MASK}{JITTER9\_SELF\_FLAT\_NO\_MASK},
      \htmlref{QUADRANT\_JITTER}{QUADRANT\_JITTER},\\
      \htmlref{QUADRANT\_JITTER\_NO\_MASK}{QUADRANT\_JITTER\_NO\_MASK}.
   }
   \sstimplementationstatus{
      \sstitemlist{

         \sstitem
         The processing engines are from the Starlink packages: \xref{\CCDPACK}{sun139}{},
         \xref{\KAPPA}{sun95}{}, \xref{\PISA}{sun109}{}, and \xref{\PHOTOM}{sun45}{}.

         \sstitem
         Uses the Starlink NDF format.

         \sstitem
         History is recorded within the data files.

         \sstitem
         The title of the data is propagated through intermediate files
         to the mosaic.

         \sstitem
         Error propagation is not used.
      }
   }
}
\newpage
\sstroutine{
   JITTER5\_SELF\_FLAT\_NO\_MASK
}{
   Reduces a 5-point ``standard jitter'' photometry observation
}{
   \sstdescription{
      This script reduces a ``standard jitter'' photometry observation
      with UFTI data.  It takes an UFTI observation comprising five
      object frames and a dark frame to make a calibrated, untrimmed
      mosaic automatically.

      It performs a null debiassing, bad-pixel masking, dark
      subtraction, flat-field creation and division, feature
      detection and matching between object frames, and resampling.

      This recipe works well for faint sources and sparse fields.
   }
   \sstnotes{
      \sstitemlist{

         \sstitem
         The bad pixel mask applied is {\tt\$ORAC\_DATA\_CAL/bpm}.

         \sstitem
         The flat field is created by combining normalised object
         frames using the clipped median at each pixel.

         \sstitem
         Registration is performed using common point sources in the
         overlap regions.  If the recipe cannot identify sufficient common
         objects, the script resorts to using the telescope offsets
         transformed to pixels.

         \sstitem
         The resampling applies non-integer shifts of origin using
         bilinear interpolation.  There is no rotation to align the
         Cartesian axes with the cardinal directions.

         \sstitem
         The recipe makes the mosaics by applying offsets in intensity
         to give the most consistent result amongst the overlapping regions.
         The mosaic is not trimmed to the dimensions of a single frame, thus
         the noise will be greater in the peripheral areas having received
         less exposure time.  The mosaic is not normalised by its exposure
         time (that being the exposure time of a single frame).

         \sstitem
         For each cycle of five, the recipe creates a mosaic, which has
         its bad pixels filled and is then added into a master mosaic of
         improving signal to noise.  The exposure time is also summed and
         stored in the mosaic's EXP\_TIME header.  Likewise the end airmass
         header, AMEND, is updated to match that of the last-observed frame
         contributing to the mosaic.

         \sstitem
         Intermediate frames are deleted except for the flat-fielded ({\tt\_ff}
         suffix) frames.
      }
   }
   \sstdiytopic{
      Output Data
   }{
      \sstitemlist{

         \sstitem
         The integrated mosaic in {\tt g$<$date$>$\_$<$group\_number$>$\_mos}.

         \sstitem
         A mosaic for each cycle of four in
         {\tt g$<$date$>$\_$<$group\_number$>$\_mos\_$<$cycle\_number$>$}, where {\tt$<$cycle\_number$>$}
         counts from 0.

         \sstitem
         The individual flat-fielded frames in {\tt f$<$date$>$\_$<$obs\_number$>$\_ff}.

         \sstitem
         The created flat fields in {\tt flat\_$<$group\_number$>$subgrp} for the
         first or only cycle, and {\tt flat\_$<$group\_number$>$subgrp\_$<$cycle\_number$>$}
         for subsequent cycles.
      }
   }
   \sstparameters{
      \sstsubsection{
         NUMBER = INTEGER
      }{
         The number of frames in the jitter pattern. {\tt[5]}
      }
   }
   \sstdiytopic{
      Related Recipes
   }{
      \htmlref{JITTER5\_SELF\_FLAT}{JITTER5\_SELF\_FLAT},
      \htmlref{JITTER5\_SELF\_FLAT\_APHOT}{JITTER5\_SELF\_FLAT\_APHOT},
      \htmlref{JITTER9\_SELF\_FLAT}{JITTER9\_SELF\_FLAT}, \\
      \htmlref{JITTER9\_SELF\_FLAT\_NO\_MASK}{JITTER9\_SELF\_FLAT\_NO\_MASK},
      \htmlref{QUADRANT\_JITTER}{QUADRANT\_JITTER},\\
      \htmlref{QUADRANT\_JITTER\_NO\_MASK}{QUADRANT\_JITTER\_NO\_MASK}.
   }
   \sstimplementationstatus{
      \sstitemlist{

         \sstitem
         The processing engines are from the Starlink packages: \xref{\CCDPACK}{sun139}{},
         \xref{\KAPPA}{sun95}{}, and  \xref{\PISA}{sun109}{}.

         \sstitem
         Uses the Starlink NDF format.

         \sstitem
         History is recorded within the data files.

         \sstitem
         The title of the data is propagated through intermediate files
         to the mosaic.

         \sstitem
         Error propagation is not used.
      }
   }
}
\newpage
\sstroutine{
   JITTER9\_SELF\_FLAT
}{
   Reduces a 9-point ``standard jitter'' photometry observation using
   object masking
}{
   \sstdescription{
      This script reduces a ``standard jitter'' photometry observation
      with UFTI data.  It takes an UFTI observation comprising nine
      object frames and a dark frame to make a calibrated, untrimmed
      mosaic automatically.

      It performs a null debiassing, bad-pixel masking, dark
      subtraction, flat-field creation and division, feature
      detection and matching between object frames, and resampling.

      This recipe works well for faint sources and for moderately
      crowded fields.
   }
   \sstnotes{
      \sstitemlist{

         \sstitem
         The bad pixel mask applied is {\tt\$ORAC\_DATA\_CAL/bpm}.

         \sstitem
         The flat field is created iteratively.  First an approximate
         flat-field is created by combining normalised object frames using
         the median at each pixel.  This flat field is applied to the object
         frames.  Sources within the flat-fielded frames are detected, and
         masked in the dark-subtracted frames.  The first stage is repeated
         but applied to the masked frames to create the final flat field.

         \sstitem
         Registration is performed using common point sources in the
         overlap regions.  If the recipe cannot identify sufficient common
         objects, the script resorts to using the telescope offsets
         transformed to pixels.

         \sstitem
         The resampling applies non-integer shifts of origin using
         bilinear interpolation.  There is no rotation to align the
         Cartesian axes with the cardinal directions.

         \sstitem
         The recipe makes the mosaics by applying offsets in intensity
         to give the most consistent result amongst the overlapping regions.
         The mosaic is not trimmed to the dimensions of a single frame, thus
         the noise will be greater in the peripheral areas having received
         less exposure time.  The mosaic is not normalised by its exposure
         time (the exposure time of a single frame).

         \sstitem
         For each cycle of nine, the recipe creates a mosaic, which has
         its bad pixels filled and is then added into a master mosaic of
         improving signal to noise.  The exposure time is also summed and
         stored in the mosaic's EXP\_TIME header.  Likewise the end airmass
         header, AMEND, is updated to match that of the last-observed frame
         contributing to the mosaic.

         \sstitem
         Intermediate frames are deleted except for the flat-fielded ({\tt\_ff}
         suffix) frames.
      }
   }
   \sstdiytopic{
      Output Data
   }{
      \sstitemlist{

         \sstitem
         The integrated mosaic in {\tt g$<$date$>$\_$<$group\_number$>$\_mos}.

         \sstitem
         A mosaic for each cycle of four in
         {\tt g$<$date$>$\_$<$group\_number$>$\_mos\_$<$cycle\_number$>$}, where {\tt$<$cycle\_number$>$}
         counts from 0.

         \sstitem
         The individual flat-fielded frames in {\tt f$<$date$>$\_$<$obs\_number$>$\_ff}.

         \sstitem
         The created flat fields in {\tt flat\_$<$group\_number$>$subgrp} for the
         first or only cycle, and {\tt flat\_$<$group\_number$>$subgrp\_$<$cycle\_number$>$}
         for subsequent cycles.
      }
   }
   \sstparameters{
      \sstsubsection{
         NUMBER = INTEGER
      }{
         The number of frames in the jitter pattern. {\tt[9]}
      }
   }
   \sstdiytopic{
      Related Recipes
   }{
      \htmlref{JITTER5\_SELF\_FLAT}{JITTER5\_SELF\_FLAT},
      \htmlref{JITTER5\_SELF\_FLAT\_NO\_MASK}{JITTER5\_SELF\_FLAT\_NO\_MASK},\\
      \htmlref{JITTER5\_SELF\_FLAT\_APHOT}{JITTER5\_SELF\_FLAT\_APHOT},
      \htmlref{JITTER9\_SELF\_FLAT\_NO\_MASK}{JITTER9\_SELF\_FLAT\_NO\_MASK}, \\
      \htmlref{QUADRANT\_JITTER}{QUADRANT\_JITTER},
      \htmlref{QUADRANT\_JITTER\_NO\_MASK}{QUADRANT\_JITTER\_NO\_MASK}.
   }
   \sstimplementationstatus{
      \sstitemlist{

         \sstitem
         The processing engines are from the Starlink packages: \xref{\CCDPACK}{sun139}{},
         \xref{\KAPPA}{sun95}{}, and  \xref{\PISA}{sun109}{}.

         \sstitem
         Uses the Starlink NDF format.

         \sstitem
         History is recorded within the data files.

         \sstitem
         The title of the data is propagated through intermediate files
         to the mosaic.

         \sstitem
         Error propagation is not used.
      }
   }
}
\newpage
\sstroutine{
   JITTER9\_SELF\_FLAT\_BASIC
}{
   Reduces a 9-point ``standard jitter'' photometry observation using
   just the basic operations for speed
}{
   \sstdescription{
      This script reduces a ``standard jitter'' photometry observation
      with UFTI data.  It takes an UFTI observation comprising nine
      object frames and a dark frame to make a calibrated, untrimmed
      mosaic automatically.

      It performs a null debiassing, bad-pixel masking, dark
      subtraction, flat-field creation, flat-field division, integer
      shifts of origin to register, and mosaicking.

      This recipe aims to keep pace with the pipeline's incoming
      data, and is intended for faint sources and for moderately
      crowded fields.
   }
   \sstnotes{
      \sstitemlist{

         \sstitem
         The bad pixel mask applied is {\tt\$ORAC\_DATA\_CAL/bpm}.

         \sstitem
         The flat field is created by combining normalised object
         frames using the clipped median at each pixel.

         \sstitem
         Registration is performed using the telescope offsets
         transformed to pixels.

         \sstitem
         There is no resampling, merely integer shifts of origin.

         \sstitem
         The recipe makes the mosaics by applying offsets in intensity
         to give the most consistent result amongst the overlapping regions.
         The mosaic is not trimmed to the dimensions of a single frame, thus
         the noise will be greater in the peripheral areas having received
         less exposure time.  The mosaic is not normalised by its exposure
         time (that being the exposure time of a single frame).

         \sstitem
         For each cycle of nine, the recipe creates a mosaic, which is
         then added into a master mosaic of improving signal to noise.  The
         exposure time is also summed and stored in the mosaic's EXP\_TIME
         header.  Likewise the end airmass header, AMEND, is updated to match
         that of the last-observed frame contributing to the mosaic.

         \sstitem
         Intermediate frames are deleted except for the flat-fielded ({\tt\_ff}
         suffix) frames.
      }
   }
   \sstdiytopic{
      Output Data
   }{
      \sstitemlist{

         \sstitem
         The integrated mosaic in {\tt g$<$date$>$\_$<$group\_number$>$\_mos}.

         \sstitem
         A mosaic for each cycle of four in
         {\tt g$<$date$>$\_$<$group\_number$>$\_mos\_$<$cycle\_number$>$}, where {\tt$<$cycle\_number$>$}
         counts from 0.

         \sstitem
         The individual flat-fielded frames in {\tt f$<$date$>$\_$<$obs\_number$>$\_ff}.

         \sstitem
         The created flat fields in {\tt flat\_$<$group\_number$>$subgrp} for the
         first or only cycle, and {\tt flat\_$<$group\_number$>$subgrp\_$<$cycle\_number$>$}
         for subsequent cycles.
      }
   }
   \sstparameters{
      \sstsubsection{
         NUMBER = INTEGER
      }{
         The number of frames in the jitter pattern. {\tt[9]}
      }
   }
   \sstdiytopic{
      Related Recipes
   }{
      \htmlref{JITTER5\_SELF\_FLAT}{JITTER5\_SELF\_FLAT},
      \htmlref{JITTER5\_SELF\_FLAT\_NO\_MASK}{JITTER5\_SELF\_FLAT\_NO\_MASK},\\
      \htmlref{JITTER5\_SELF\_FLAT\_APHOT}{JITTER5\_SELF\_FLAT\_APHOT},
      \htmlref{JITTER9\_SELF\_FLAT}{JITTER9\_SELF\_FLAT}, \\
      \htmlref{JITTER9\_SELF\_FLAT\_NO\_MASK}{JITTER9\_SELF\_FLAT\_NO\_MASK},
      \htmlref{QUADRANT\_JITTER}{QUADRANT\_JITTER}, \\
      \htmlref{QUADRANT\_JITTER\_NO\_MASK}{QUADRANT\_JITTER\_NO\_MASK}.
   }
   \sstimplementationstatus{
      \sstitemlist{

         \sstitem
         The processing engines are from the Starlink packages: \xref{\CCDPACK}{sun139}{},
         \xref{\KAPPA}{sun95}{}, and  \xref{\PISA}{sun109}{}.

         \sstitem
         Uses the Starlink NDF format.

         \sstitem
         History is recorded within the data files.

         \sstitem
         The title of the data is propagated through intermediate files
         to the mosaic.

         \sstitem
         Error propagation is not used.
      }
   }
}
\newpage
\sstroutine{
   JITTER9\_SELF\_FLAT\_NO\_MASK
}{
   Reduces a 9-point ``standard jitter'' photometry observation
}{
   \sstdescription{
      This script reduces a ``standard jitter'' photometry observation
      with UFTI data.  It takes an UFTI observation comprising nine
      object frames and a dark frame to make a calibrated, untrimmed
      mosaic automatically.

      It performs a null debiassing, bad-pixel masking, dark
      subtraction, flat-field creation and division, feature
      detection and matching between object frames, and resampling.

      This recipe works well for faint sources and sparse fields.
   }
   \sstnotes{
      \sstitemlist{

         \sstitem
         The bad pixel mask applied is {\tt\$ORAC\_DATA\_CAL/bpm}.

         \sstitem
         The flat field is created by combining normalised object
         frames using the clipped median at each pixel.

         \sstitem
         Registration is performed using common point sources in the
         overlap regions.  If the recipe cannot identify sufficient common
         objects, the script resorts to using the telescope offsets
         transformed to pixels.

         \sstitem
         The resampling applies non-integer shifts of origin using
         bilinear interpolation.  There is no rotation to align the
         Cartesian axes with the cardinal directions.

         \sstitem
         The recipe makes the mosaics by applying offsets in intensity
         to give the most consistent result amongst the overlapping regions.
         The mosaic is not trimmed to the dimensions of a single frame, thus
         the noise will be greater in the peripheral areas having received
         less exposure time.  The mosaic is not normalised by its exposure
         time (that being the exposure time of a single frame).

         \sstitem
         For each cycle of nine, the recipe creates a mosaic, which has
         its bad pixels filled and is then added into a master mosaic of
         improving signal to noise.  The exposure time is also summed and
         stored in the mosaic's EXP\_TIME header.  Likewise the end airmass
         header, AMEND, is updated to match that of the last-observed frame
         contributing to the mosaic.

         \sstitem
         Intermediate frames are deleted except for the flat-fielded ({\tt\_ff}
         suffix) frames.
      }
   }
   \sstdiytopic{
      Output Data
   }{
      \sstitemlist{

         \sstitem
         The integrated mosaic in {\tt g$<$date$>$\_$<$group\_number$>$\_mos}.

         \sstitem
         A mosaic for each cycle of four in
         {\tt g$<$date$>$\_$<$group\_number$>$\_mos\_$<$cycle\_number$>$}, where {\tt$<$cycle\_number$>$}
         counts from 0.

         \sstitem
         The individual flat-fielded frames in {\tt f$<$date$>$\_$<$obs\_number$>$\_ff}.

         \sstitem
         The created flat fields in {\tt flat\_$<$group\_number$>$subgrp} for the
         first or only cycle, and {\tt flat\_$<$group\_number$>$subgrp\_$<$cycle\_number$>$}
         for subsequent cycles.
      }
   }
   \sstparameters{
      \sstsubsection{
         NUMBER = INTEGER
      }{
         The number of frames in the jitter pattern. {\tt[9]}
      }
   }
   \sstdiytopic{
      Related Recipes
   }{
      \htmlref{JITTER5\_SELF\_FLAT}{JITTER5\_SELF\_FLAT},
      \htmlref{JITTER5\_SELF\_FLAT\_APHOT}{JITTER5\_SELF\_FLAT\_APHOT},\\
      \htmlref{JITTER5\_SELF\_FLAT\_NO\_MASK}{JITTER5\_SELF\_FLAT\_NO\_MASK},
      \htmlref{JITTER9\_SELF\_FLAT}{JITTER9\_SELF\_FLAT},\\
      \htmlref{QUADRANT\_JITTER}{QUADRANT\_JITTER},
      \htmlref{QUADRANT\_JITTER\_NO\_MASK}{QUADRANT\_JITTER\_NO\_MASK}.
   }
   \sstimplementationstatus{
      \sstitemlist{

         \sstitem
         The processing engines are from the Starlink packages: \xref{\CCDPACK}{sun139}{},
         \xref{\KAPPA}{sun95}{}, and  \xref{\PISA}{sun109}{}.

         \sstitem
         Uses the Starlink NDF format.

         \sstitem
         History is recorded within the data files.

         \sstitem
         The title of the data is propagated through intermediate files
         to the mosaic.

         \sstitem
         Error propagation is not used.
      }
   }
}
\newpage
\sstroutine{
   JITTER9\_SELF\_FLAT\_TELE
}{
   Reduces a 9-point ``standard jitter'' photometry observation using
   object masking, and telescope offsets for registration
}{
   \sstdescription{
      This script reduces a ``standard jitter'' photometry observation
      with UFTI data.  It takes an observation comprising nine
      object frames and a dark frame to make a calibrated, untrimmed
      mosaic automatically.

      It performs a null debiassing, bad-pixel masking, dark
      subtraction, flat-field creation and division, registration
      using telescope offsets, and resampling.

      This recipe works well for faint sources and for moderately
      crowded fields.  It is also using for observations that track a
      moving object.
   }
   \sstnotes{
      \sstitemlist{

         \sstitem
         The bad pixel mask applied is {\tt\$ORAC\_DATA\_CAL/bpm}.

         \sstitem
         The flat field is created iteratively.  First an approximate
         flat-field is created by combining normalised object frames using
         the median at each pixel.  This flat field is applied to the object
         frames.  Sources within the flat-fielded frames are detected, and
         masked in the dark-subtracted frames.  The first stage is repeated
         but applied to the masked frames to create the final flat field.

         \sstitem
         Registration is performed using the telescope offsets
         transformed to pixels.

         \sstitem
         The resampling applies non-integer shifts of origin using
         bilinear interpolation.  There is no rotation to align the
         Cartesian axes with the cardinal directions.

         \sstitem
         The recipe makes the mosaics by applying offsets in intensity
         to give the most consistent result amongst the overlapping regions.
         The mosaic is not trimmed to the dimensions of a single frame, thus
         the noise will be greater in the peripheral areas having received
         less exposure time.  The mosaic is not normalised by its exposure
         time (that being the exposure time of a single frame).

         \sstitem
         For each cycle of nine, the recipe creates a mosaic, which has
         its bad pixels filled and is then added into a master mosaic of
         improving signal to noise.  The exposure time is also summed and
         stored in the mosaic's EXP\_TIME header.  Likewise the end airmass
         header, AMEND, is updated to match that of the last-observed frame
         contributing to the mosaic.

         \sstitem
         Intermediate frames are deleted except for the flat-fielded ({\tt\_ff}
         suffix) frames.
      }
   }
   \sstdiytopic{
      Output Data
   }{
      \sstitemlist{

         \sstitem
         The integrated mosaic in {\tt g$<$date$>$\_$<$group\_number$>$\_mos}.

         \sstitem
         A mosaic for each cycle of four in
         {\tt g$<$date$>$\_$<$group\_number$>$\_mos\_$<$cycle\_number$>$}, where {\tt$<$cycle\_number$>$}
         counts from 0.

         \sstitem
         The individual flat-fielded frames in {\tt f$<$date$>$\_$<$obs\_number$>$\_ff}.

         \sstitem
         The created flat fields in {\tt flat\_$<$group\_number$>$subgrp} for the
         first or only cycle, and {\tt flat\_$<$group\_number$>$subgrp\_$<$cycle\_number$>$}
         for subsequent cycles.
      }
   }
   \sstparameters{
      \sstsubsection{
         NUMBER = INTEGER
      }{
         The number of frames in the jitter pattern. {\tt[9]}
      }
   }
   \sstdiytopic{
      Related Recipes
   }{
      \htmlref{JITTER5\_SELF\_FLAT}{JITTER5\_SELF\_FLAT},
      \htmlref{JITTER5\_SELF\_FLAT\_NO\_MASK}{JITTER5\_SELF\_FLAT\_NO\_MASK},\\
      \htmlref{JITTER5\_SELF\_FLAT\_APHOT}{JITTER5\_SELF\_FLAT\_APHOT},
      \htmlref{JITTER9\_SELF\_FLAT\_NO\_MASK}{JITTER9\_SELF\_FLAT\_NO\_MASK},\\
      \htmlref{QUADRANT\_JITTER}{QUADRANT\_JITTER},
      \htmlref{QUADRANT\_JITTER\_NO\_MASK}{QUADRANT\_JITTER\_NO\_MASK}.
   }
   \sstimplementationstatus{
      \sstitemlist{

         \sstitem
         The processing engines are from the Starlink packages: \xref{\CCDPACK}{sun139}{},
         \xref{\KAPPA}{sun95}{}, and  \xref{\PISA}{sun109}{}.

         \sstitem
         Uses the Starlink NDF format.

         \sstitem
         History is recorded within the data files.

         \sstitem
         The title of the data is propagated through intermediate files
         to the mosaic.

         \sstitem
         Error propagation is not used.
      }
   }
}
\newpage
\sstroutine{
   JITTER\_SELF\_FLAT
}{
   Reduces a 9-point ``standard jitter'' photometry observation using
   object masking
}{
   \sstdescription{
      This script reduces a ``standard jitter'' photometry observation
      with UFTI data.  It takes an UFTI observation comprising nine
      object frames and a dark frame to make a calibrated, untrimmed
      mosaic automatically.

      It performs a null debiassing, bad-pixel masking, dark
      subtraction, flat-field creation and division, feature
      detection and matching between object frames, and resampling.

      This recipe works well for faint sources and for moderately
      crowded fields.
   }
   \sstnotes{
      \sstitemlist{

         \sstitem
         The bad pixel mask applied is {\tt\$ORAC\_DATA\_CAL/bpm}.

         \sstitem
         The flat field is created iteratively.  First an approximate
         flat-field is created by combining normalised object frames using
         the median at each pixel.  This flat field is applied to the object
         frames.  Sources within the flat-fielded frames are detected, and
         masked in the dark-subtracted frames.  The first stage is repeated
         but applied to the masked frames to create the final flat field.

         \sstitem
         Registration is performed using common point sources in the
         overlap regions.  If the recipe cannot identify sufficient common
         objects, the script resorts to using the telescope offsets
         transformed to pixels.

         \sstitem
         The resampling applies non-integer shifts of origin using
         bilinear interpolation.  There is no rotation to align the
         Cartesian axes with the cardinal directions.

         \sstitem
         The recipe makes the mosaics by applying offsets in intensity
         to give the most consistent result amongst the overlapping regions.
         The mosaic is not trimmed to the dimensions of a single frame, thus
         the noise will be greater in the peripheral areas having received
         less exposure time.  The mosaic is not normalised by its exposure
         time (that being the exposure time of a single frame).

         \sstitem
         For each cycle of nine, the recipe creates a mosaic, which has
         its bad pixels filled and is then added into a master mosaic of
         improving signal to noise.  The exposure time is also summed and
         stored in the mosaic's EXP\_TIME header.  Likewise the end airmass
         header, AMEND, is updated to match that of the last-observed frame
         contributing to the mosaic.

         \sstitem
         Intermediate frames are deleted except for the flat-fielded ({\tt\_ff}
         suffix) frames.
      }
   }
   \sstdiytopic{
      Output Data
   }{
      \sstitemlist{

         \sstitem
         The integrated mosaic in {\tt g$<$date$>$\_$<$group\_number$>$\_mos}.

         \sstitem
         A mosaic for each cycle of four in
         {\tt g$<$date$>$\_$<$group\_number$>$\_mos\_$<$cycle\_number$>$}, where {\tt$<$cycle\_number$>$}
         counts from 0.

         \sstitem
         The individual flat-fielded frames in {\tt f$<$date$>$\_$<$obs\_number$>$\_ff}.

         \sstitem
         The created flat fields in {\tt flat\_$<$group\_number$>$subgrp} for the
         first or only cycle, and {\tt flat\_$<$group\_number$>$subgrp\_$<$cycle\_number$>$}
         for subsequent cycles.
      }
   }
   \sstparameters{
      \sstsubsection{
         NUMBER = INTEGER
      }{
         The number of frames in the jitter pattern. {\tt[9]}
      }
   }
   \sstdiytopic{
      Related Recipes
   }{
      \htmlref{JITTER5\_SELF\_FLAT}{JITTER5\_SELF\_FLAT},
      \htmlref{JITTER5\_SELF\_FLAT\_NO\_MASK}{JITTER5\_SELF\_FLAT\_NO\_MASK},
      \htmlref{JITTER9\_SELF\_FLAT}{JITTER9\_SELF\_FLAT}, \\
      \htmlref{JITTER9\_SELF\_FLAT\_NO\_MASK}{JITTER9\_SELF\_FLAT\_NO\_MASK},
      \htmlref{QUADRANT\_JITTER}{QUADRANT\_JITTER},\\
      \htmlref{QUADRANT\_JITTER\_NO\_MASK}{QUADRANT\_JITTER\_NO\_MASK}.
   }
   \sstimplementationstatus{
      \sstitemlist{

         \sstitem
         The processing engines are from the Starlink packages: \xref{\CCDPACK}{sun139}{},
         \xref{\KAPPA}{sun95}{}, and  \xref{\PISA}{sun109}{}.

         \sstitem
         Uses the Starlink NDF format.

         \sstitem
         History is recorded within the data files.

         \sstitem
         The title of the data is propagated through intermediate files
         to the mosaic.

         \sstitem
         Error propagation is not used.
      }
   }
}
\newpage
\sstroutine{
   MOVING\_JITTER9\_SELF\_FLAT
}{
   Reduces a 9-point ``standard jitter'' photometry observation of a
   moving target using object masking
}{
   \sstdescription{
      This script reduces a ``standard jitter'' photometry observation
      with UFTI data.  It takes an observation comprising nine
      object frames and a dark frame to make a calibrated, untrimmed
      mosaic automatically.

      It performs a null debiassing, bad-pixel masking, dark
      subtraction, flat-field creation and division, feature
      detection and matching between object frames, and resampling.
      Registration is adjusted to track the motion of the moving
      target using ephemeris data stored in file {\tt target\_ephem.txt}.
      See \htmlref{``Ephemeris-file Format''}{ephem_format} for details
      of this file's format.

      This recipe works well for faint moving sources and in moderately
      crowded fields.  It should not be used for frames where the
      telescope guided on the moving object.  In that case reduction
      should be performed by \htmlref{JITTER9\_SELF\_FLAT\_TELE}{JITTER9\_SELF\_FLAT\_TELE} which registers
      using the telescope offsets alone.
   }
   \sstnotes{
      \sstitemlist{

         \sstitem
         The bad pixel mask applied is {\tt\$ORAC\_DATA\_CAL/bpm}.

         \sstitem
         The flat field is created iteratively.  First an approximate
         flat-field is created by combining normalised object frames using
         the median at each pixel.  This flat field is applied to the object
         frames.  Sources within the flat-fielded frames are detected, and
         masked in the dark-subtracted frames.  The first stage is repeated
         but applied to the masked frames to create the final flat field.

         \sstitem
         Registration is performed using common point sources in the
         overlap regions.  If the recipe cannot identify sufficient common
         objects, the script resorts to using the telescope offsets
         transformed to pixels.  Once the offsets are determined, they
         are adjusted for the motion of the target, so that the final
         mosaic registers the target, not the background stars.

         \sstitem
         The ephemeris file is specified by environment variable
         ORAC\_EPHEMERIS, defaulting to {\tt \$ORAC\_DATA\_OUT/target\_ephem.dat}.

         \sstitem
         The resampling applies non-integer shifts of origin using
         bilinear interpolation.  There is no rotation to align the
         Cartesian axes with the cardinal directions.

         \sstitem
         The recipe makes the mosaics by applying offsets in intensity
         to give the most consistent result amongst the overlapping regions.
         The mosaic is not trimmed to the dimensions of a single frame, thus
         the noise will be greater in the peripheral areas having received
         less exposure time.  The mosaic is not normalised by its exposure
         time (that being the exposure time of a single frame).

         \sstitem
         For each cycle of nine, the recipe creates a mosaic, which has
         its bad pixels filled and is then added into a master mosaic of
         improving signal to noise.  The exposure time is also summed and
         stored in the mosaic's EXP\_TIME header.  Likewise the end airmass
         header, AMEND, is updated to match that of the last-observed frame
         contributing to the mosaic.

         \sstitem
         Intermediate frames are deleted except for the flat-fielded ({\tt\_ff}
         suffix) frames.
      }
   }
   \label{ephem_format}
   \sstdiytopic{
      Ephemeris-file Format
   }{
      The current format of the ephemeris file is one line per object
      comprising three space-separated fields in the following order:
      \ssthitemlist{

         \sstitem
           the objectname, which may contain embedded spaces;

         \sstitem
           the motion in the plane of the sky in arcseconds/second for right
           ascension then declination.

      }
      Note that the right ascension motion is the change in right ascension
      multiplied by the cosine of the declination.  The format will change
      to include UT and possibly date.
   }
   \sstdiytopic{
      Output Data
   }{
      \sstitemlist{

         \sstitem
         The integrated mosaic in {\tt g$<$date$>$\_$<$group\_number$>$\_mos}.

         \sstitem
         A mosaic for each cycle of four in
         {\tt g$<$date$>$\_$<$group\_number$>$\_mos\_$<$cycle\_number$>$}, where {\tt$<$cycle\_number$>$}
         counts from 0.

         \sstitem
         The individual flat-fielded frames in {\tt f$<$date$>$\_$<$obs\_number$>$\_ff}.

         \sstitem
         The created flat fields in {\tt flat\_$<$group\_number$>$subgrp} for the
         first or only cycle, and {\tt flat\_$<$group\_number$>$subgrp\_$<$cycle\_number$>$}
         for subsequent cycles.
      }
   }
   \sstparameters{
      \sstsubsection{
         NUMBER = INTEGER
      }{
         The number of frames in the jitter pattern. {\tt[9]}
      }
   }
   \sstdiytopic{
      Related Recipes
   }{
      \htmlref{JITTER9\_SELF\_FLAT}{JITTER9\_SELF\_FLAT},
      \htmlref{JITTER9\_SELF\_FLAT\_TELE}{JITTER9\_SELF\_FLAT\_TELE},\\
      \htmlref{MOVING\_JITTER9\_SELF\_FLAT\_BASIC}{MOVING\_JITTER9\_SELF\_FLAT\_BASIC}.
   }
   \sstimplementationstatus{
      \sstitemlist{

         \sstitem
         The processing engines are from the Starlink packages: \xref{\CCDPACK}{sun139}{},
         \xref{\KAPPA}{sun95}{}, and  \xref{\PISA}{sun109}{}.

         \sstitem
         Uses the Starlink NDF format.

         \sstitem
         History is recorded within the data files.

         \sstitem
         The title of the data is propagated through intermediate files
         to the mosaic.

         \sstitem
         Error propagation is not used.
      }
   }
}
\newpage
\sstroutine{
   MOVING\_JITTER9\_SELF\_FLAT\_BASIC
}{
   Reduces a 9-point ``standard jitter'' photometry observation of a
   moving target using just the basic operations for speed
}{
   \sstdescription{
      This script reduces a ``standard jitter'' photometry observation
      with UFTI data.  It takes an observation comprising nine
      object frames and a dark frame to make a calibrated, untrimmed
      mosaic automatically.

      It performs a null debiassing, bad-pixel masking, dark
      subtraction, flat-field creation and division, amd integer shifts
      of pixel origin to register to fixed sky co-ordinates.  The
      registration is adjusted to track the motion of the moving
      target using ephemeris data stored in file {\tt target\_ephem.txt}.
      See \htmlref{``Ephemeris-file Format''}{ephem_format} for details
      of this file's format.

      This recipe aims to keep pace with the pipeline's incoming data.
      It works well for faint moving sources and in moderately
      crowded fields.  It should not be used for frames where the
      telescope guided on the moving object.  In that case reduction
      should be performed by \htmlref{JITTER9\_SELF\_FLAT\_TELE}{JITTER9\_SELF\_FLAT\_TELE} which registers
      using the telescope offsets alone.
   }
   \sstnotes{
      \sstitemlist{

         \sstitem
         The bad pixel mask applied is {\tt\$ORAC\_DATA\_CAL/bpm}.

         \sstitem
         The flat field is created by combining normalised object
         frames using the median at each pixel.

         \sstitem
         Registration is performed using the telescope offsets
         transformed to pixels.  Once the offsets are determined, they
         are adjusted for the motion of the target, so that the final
         mosaic registers the target, not the background stars.

         \sstitem
         There is no resampling, merely integer shifts of origin.

         \sstitem
         The ephemeris file is specified by environment variable
         ORAC\_EPHEMERIS, defaulting to {\tt \$ORAC\_DATA\_OUT/target\_ephem.dat}.

         \sstitem
         The recipe makes the mosaics by applying offsets in intensity
         to give the most consistent result amongst the overlapping regions.
         The mosaic is not trimmed to the dimensions of a single frame, thus
         the noise will be greater in the peripheral areas having received
         less exposure time.  The mosaic is not normalised by its exposure
         time (that being the exposure time of a single frame).

         \sstitem
         For each cycle of nine, the recipe creates a mosaic, which is
         then added into a master mosaic of improving signal to noise.  The
         exposure time is also summed and stored in the mosaic's EXP\_TIME
         header.  Likewise the end airmass header, AMEND, is updated to match
         that of the last-observed frame contributing to the mosaic.

         \sstitem
         Intermediate frames are deleted except for the flat-fielded ({\tt\_ff}
         suffix) frames.
      }
   }
   \sstdiytopic{
      Ephemeris-file Format
   }{
       The current format of the ephemeris file is one line per object
       comprising three space-separated fields in the following order:
      \ssthitemlist{

         \sstitem
            the objectname, which may contain embedded spaces; and

         \sstitem
            the motion in the plane of the sky in arcseconds/second for right
          ascension then declination.

      }
      Note that the right-ascension motion is the change in right ascension
      multiplied by the cosine of the declination.  The format may change
      to include UT and possibly date.
   }
   \sstdiytopic{
      Output Data
   }{
      \sstitemlist{

         \sstitem
         The integrated mosaic in {\tt g$<$date$>$\_$<$group\_number$>$\_mos}.

         \sstitem
         A mosaic for each cycle of four in
         {\tt g$<$date$>$\_$<$group\_number$>$\_mos\_$<$cycle\_number$>$}, where {\tt$<$cycle\_number$>$}
         counts from 0.

         \sstitem
         The individual flat-fielded frames in {\tt f$<$date$>$\_$<$obs\_number$>$\_ff}.

         \sstitem
         The created flat fields in {\tt flat\_$<$group\_number$>$subgrp} for the
         first or only cycle, and {\tt flat\_$<$group\_number$>$subgrp\_$<$cycle\_number$>$}
         for subsequent cycles.
      }
   }
   \sstparameters{
      \sstsubsection{
         NUMBER = INTEGER
      }{
         The number of frames in the jitter pattern. {\tt[9]}
      }
   }
   \sstdiytopic{
      Related Recipes
   }{
      \htmlref{JITTER9\_SELF\_FLAT}{JITTER9\_SELF\_FLAT},
      \htmlref{JITTER9\_SELF\_FLAT\_BASIC}{JITTER9\_SELF\_FLAT\_BASIC},\\
      \htmlref{JITTER9\_SELF\_FLAT\_TELE}{JITTER9\_SELF\_FLAT\_TELE},
      \htmlref{MOVING\_JITTER9\_SELF\_FLAT}{MOVING\_JITTER9\_SELF\_FLAT}.
   }
   \sstimplementationstatus{
      \sstitemlist{

         \sstitem
         The processing engines are from the Starlink packages: \xref{\CCDPACK}{sun139}{}
         and \xref{\KAPPA}{sun95}{}.

         \sstitem
         Uses the Starlink NDF format.

         \sstitem
         History is recorded within the data files.

         \sstitem
         The title of the data is propagated through intermediate files
         to the mosaic.

         \sstitem
         Error propagation is not used.
      }
   }
}
\newpage
\sstroutine{
   NIGHT\_LOG
}{
   Produces a text file log of a night's observations
}{
   \sstdescription{
      This recipe takes a night's observations, and creates a text file
      containing a headed tabulation of parameters for each frame.

      The parameters are: observation number, object name, observation type,
      UT start time, exposure time, number of coadds, read mode and speed,
      filter, start airmass, frame dimensions in pixels, base equatorial
      co-ordinates, and data-reduction recipe name.
   }
   \sstnotes{
      \sstitemlist{

         \sstitem
         Run with {\tt oracdr -noeng -from 1 -skip}  for efficiency.

         \sstitem
         The {\tt $<$date$>$} comes from the header keyword DATE.

         \sstitem
         Specification provided by Sandy Leggett.
      }
   }
   \sstdiytopic{
      Output Data
   }{
      \sstitemlist{

         \sstitem
         The text log file {\tt \$ORAC\_DATA\_IN/$<$date$>$.nightlog}, where
         {\tt $<$date$>$} is the UT date.
      }
   }
}
\newpage
\sstroutine{
   NOD4\_SELF\_FLAT\_NO\_MASK
}{
   Reduces a 4-point ``nod jitter'' observation
}{
   \sstdescription{
      This script reduces a ``nod jitter'' observation with UFTI data.  It
      takes an UFTI observation comprising four object frames and a dark
      frame to make a calibrated, untrimmed mosaic automatically.

      It performs a null debiassing, bad-pixel masking, dark
      subtraction, flat-field creation and division, feature
      detection and matching between object frames, and resampling.

      This recipe works well for faint sources in moderately crowded fields.
   }
   \sstnotes{
      \sstitemlist{

         \sstitem
         The bad pixel mask applied is {\tt\$ORAC\_DATA\_CAL/bpm}.

         \sstitem
         The flat field is created by combining normalised object
         frames using the median at each pixel.

         \sstitem
         Registration is performed using common point sources in the
         overlap regions.  If the recipe cannot identify sufficient common
         objects, the script resorts to using the telescope offsets
         transformed to pixels.

         \sstitem
         The resampling applies non-integer shifts of origin using
         bilinear interpolation.  There is no rotation to align the
         Cartesian axes with the cardinal directions.

         \sstitem
         The recipe makes the mosaics by applying offsets in intensity
         to give the most consistent result amongst the overlapping regions.
         The mosaic is not trimmed to the dimensions of a single frame, thus
         the noise will be greater in the peripheral areas having received
         less exposure time.  The mosaic is not normalised by its exposure
         time (that being the exposure time of a single frame).

         \sstitem
         For each cycle of four, the recipe creates a mosaic, which has
         its bad pixels filled and is then added into a master mosaic of
         improving signal to noise.  The exposure time is also summed and
         stored in the mosaic's EXP\_TIME header.  Likewise the end airmass
         header, AMEND, is updated to match that of the last-observed frame
         contributing to the mosaic.

         \sstitem
         Intermediate frames are deleted except for the flat-fielded ({\tt\_ff}
         suffix) frames.
      }
   }
   \sstdiytopic{
      Output Data
   }{
      \sstitemlist{

         \sstitem
         The integrated mosaic in {\tt g$<$date$>$\_$<$group\_number$>$\_mos}.

         \sstitem
         A mosaic for each cycle of four in
         {\tt g$<$date$>$\_$<$group\_number$>$\_mos\_$<$cycle\_number$>$}, where {\tt$<$cycle\_number$>$}
         counts from 0.

         \sstitem
         The individual flat-fielded frames in {\tt f$<$date$>$\_$<$obs\_number$>$\_ff}.

         \sstitem
         The created flat fields in {\tt flat\_$<$group\_number$>$subgrp} for the
         first or only cycle, and {\tt flat\_$<$group\_number$>$subgrp\_$<$cycle\_number$>$}
         for subsequent cycles.
      }
   }
   \sstparameters{
      \sstsubsection{
         NUMBER = INTEGER
      }{
         The number of frames in the nod pattern.  It must be a multiple
         of four. {\tt[4]}
      }
   }
   \sstdiytopic{
      Related Recipes
   }{
      \htmlref{NOD4\_SELF\_FLAT\_NO\_MASK\_APHOT}{NOD4\_SELF\_FLAT\_NO\_MASK\_APHOT},
      \htmlref{NOD8\_SELF\_FLAT\_NO\_MASK}{NOD8\_SELF\_FLAT\_NO\_MASK},\\
      \htmlref{NOD8\_SELF\_FLAT\_NO\_MASK\_APHOT}{NOD8\_SELF\_FLAT\_NO\_MASK\_APHOT}.
   }
   \sstimplementationstatus{
      \sstitemlist{

         \sstitem
         The processing engines are from the Starlink packages: \xref{\CCDPACK}{sun139}{},
         \xref{\KAPPA}{sun95}{}, and  \xref{\PISA}{sun109}{}.

         \sstitem
         Uses the Starlink NDF format.

         \sstitem
         History is recorded within the data files.

         \sstitem
         The title of the data is propagated through intermediate files
         to the mosaic.

         \sstitem
         Error propagation is not used.
      }
   }
}
\newpage
\sstroutine{
   NOD4\_SELF\_FLAT\_NO\_MASK\_APHOT
}{
   Reduces a 4-point ``nod jitter'' photometry observation,
   and performs aperture photometry
}{
   \sstdescription{
      This script reduces a ``nod jitter'' photometry observation with
      UFTI data.  It takes an UFTI observation comprising four
      object frames and a dark frame to make a calibrated, untrimmed
      mosaic automatically.

      It performs a null debiassing, bad-pixel masking, dark
      subtraction, flat-field creation and division, feature
      detection and matching between object frames, and resampling.

      Photometry of the point source using a fixed 5-arcsecond aperture
      is calculated for each jitter frame and the mosaic.  The results
      appear in {\tt\$ORAC\_DATA\_OUT/aphot\_results.txt} in the form of a Starlink
      \xref{small text list}{sun190}{STLREF}.  The analysis of each star is appended to this file.

      This recipe works well for faint sources in moderately crowded fields.
   }
   \sstnotes{
      \sstitemlist{

         \sstitem
         The bad pixel mask applied is {\tt\$ORAC\_DATA\_CAL/bpm}.

         \sstitem
         The flat field is created by combining normalised object
         frames using the median at each pixel.

         \sstitem
         Registration is performed using common point sources in the
         overlap regions.  If the recipe cannot identify sufficient common
         objects, the script resorts to using the telescope offsets
         transformed to pixels.

         \sstitem
         The resampling applies non-integer shifts of origin using
         bilinear interpolation.  There is no rotation to align the
         Cartesian axes with the cardinal directions.

         \sstitem
         The recipe makes the mosaics by applying offsets in intensity
         to give the most consistent result amongst the overlapping regions.
         The mosaic is not trimmed to the dimensions of a single frame, thus
         the noise will be greater in the peripheral areas having received
         less exposure time.  The mosaic is not normalised by its exposure
         time (that being the exposure time of a single frame).

         \sstitem
         For each cycle of four, the recipe creates a mosaic, which has
         its bad pixels filled and is then added into a master mosaic of
         improving signal to noise.  The exposure time is also summed and
         stored in the mosaic's EXP\_TIME header.  Likewise the end airmass
         header, AMEND, is updated to match that of the last-observed frame
         contributing to the mosaic.

         \sstitem
         Intermediate frames are deleted except for the flat-fielded ({\tt\_ff}
         suffix) frames.
      }
   }
   \sstdiytopic{
      Output Data
   }{
      \sstitemlist{

         \sstitem
         The integrated mosaic in {\tt g$<$date$>$\_$<$group\_number$>$\_mos}.

         \sstitem
         A mosaic for each cycle of four in
         {\tt g$<$date$>$\_$<$group\_number$>$\_mos\_$<$cycle\_number$>$}, where {\tt$<$cycle\_number$>$}
         counts from 0.

         \sstitem
         The individual flat-fielded frames in {\tt f$<$date$>$\_$<$obs\_number$>$\_ff}.

         \sstitem
         The created flat fields in {\tt flat\_$<$group\_number$>$subgrp} for the
         first or only cycle, and {\tt flat\_$<$group\_number$>$subgrp\_$<$cycle\_number$>$}
         for subsequent cycles.
      }
   }
   \sstparameters{
      \sstsubsection{
         NUMBER = INTEGER
      }{
         The number of frames in the nod pattern.  It must be a multiple
         of four. {\tt[4]}
      }
   }
   \sstdiytopic{
      Related Recipes
   }{
      \htmlref{NOD8\_SELF\_FLAT\_NO\_MASK\_APHOT}{NOD8\_SELF\_FLAT\_NO\_MASK\_APHOT}.
   }
   \sstimplementationstatus{
      \sstitemlist{

         \sstitem
         The processing engines are from the Starlink packages: \xref{\CCDPACK}{sun139}{},
         \xref{\KAPPA}{sun95}{}, \xref{\PISA}{sun109}{}, and \xref{\PHOTOM}{sun45}{}.

         \sstitem
         Uses the Starlink NDF format.

         \sstitem
         History is recorded within the data files.

         \sstitem
         The title of the data is propagated through intermediate files
         to the mosaic.

         \sstitem
         Error propagation is not used.
      }
   }
}
\newpage
\sstroutine{
   NOD8\_SELF\_FLAT\_NO\_MASK
}{
   Reduces a 8-point ``nod jitter'' observation
}{
   \sstdescription{
      This script reduces a ``nod jitter'' observation with UFTI data.
      It takes an UFTI observation comprising eight object frames and a
      dark frame to make a calibrated, untrimmed  mosaic automatically.

      It performs a null debiassing, bad-pixel masking, dark
      subtraction, flat-field creation and division, feature
      detection and matching between object frames, and resampling.

      This recipe works well for faint sources in moderately crowded fields.
   }
   \sstnotes{
      \sstitemlist{

         \sstitem
         The bad pixel mask applied is {\tt\$ORAC\_DATA\_CAL/bpm}.

         \sstitem
         The flat field is created by combining normalised object
         frames using the median at each pixel.

         \sstitem
         Registration is performed using common point sources in the
         overlap regions.  If the recipe cannot identify sufficient common
         objects, the script resorts to using the telescope offsets
         transformed to pixels.

         \sstitem
         The resampling applies non-integer shifts of origin using
         bilinear interpolation.  There is no rotation to align the
         Cartesian axes with the cardinal directions.

         \sstitem
         The recipe makes the mosaics by applying offsets in intensity
         to give the most consistent result amongst the overlapping regions.
         The mosaic is not trimmed to the dimensions of a single frame, thus
         the noise will be greater in the peripheral areas having received
         less exposure time.  The mosaic is not normalised by its exposure
         time (that being the exposure time of a single frame).

         \sstitem
         For each cycle of eight, the recipe creates a mosaic, which has
         its bad pixels filled and is then added into a master mosaic of
         improving signal to noise.  The exposure time is also summed and
         stored in the mosaic's EXP\_TIME header.  Likewise the end airmass
         header, AMEND, is updated to match that of the last-observed frame
         contributing to the mosaic.

         \sstitem
         Intermediate frames are deleted except for the flat-fielded ({\tt\_ff}
         suffix) frames.
      }
   }
   \sstdiytopic{
      Output Data
   }{
      \sstitemlist{

         \sstitem
         The integrated mosaic in {\tt g$<$date$>$\_$<$group\_number$>$\_mos}.

         \sstitem
         A mosaic for each cycle of eight in
         {\tt g$<$date$>$\_$<$group\_number$>$\_mos\_$<$cycle\_number$>$}, where {\tt$<$cycle\_number$>$}
         counts from 0.

         \sstitem
         The individual flat-fielded frames in {\tt f$<$date$>$\_$<$obs\_number$>$\_ff}.

         \sstitem
         The created flat fields in {\tt flat\_$<$group\_number$>$subgrp} for the
         first or only cycle, and {\tt flat\_$<$group\_number$>$subgrp\_$<$cycle\_number$>$}
         for subsequent cycles.
      }
   }
   \sstparameters{
      \sstsubsection{
         NUMBER = INTEGER
      }{
         The number of frames in the nod pattern.  It must be a multiple
         of four. {\tt[8]}
      }
   }
   \sstdiytopic{
      Related Recipes
   }{
      \htmlref{NOD4\_SELF\_FLAT\_NO\_MASK\_APHOT}{NOD4\_SELF\_FLAT\_NO\_MASK\_APHOT},
      \htmlref{NOD4\_SELF\_FLAT\_NO\_MASK}{NOD4\_SELF\_FLAT\_NO\_MASK},\\
      \htmlref{NOD8\_SELF\_FLAT\_NO\_MASK\_APHOT}{NOD8\_SELF\_FLAT\_NO\_MASK\_APHOT}.
   }
   \sstimplementationstatus{
      \sstitemlist{

         \sstitem
         The processing engines are from the Starlink packages: \xref{\CCDPACK}{sun139}{},
         \xref{\KAPPA}{sun95}{}, and  \xref{\PISA}{sun109}{}.

         \sstitem
         Uses the Starlink NDF format.

         \sstitem
         History is recorded within the data files.

         \sstitem
         The title of the data is propagated through intermediate files
         to the mosaic.

         \sstitem
         Error propagation is not used.
      }
   }
}
\newpage
\sstroutine{
   NOD8\_SELF\_FLAT\_NO\_MASK\_APHOT
}{
   Reduces a 8-point ``nod jitter'' photometry observation,
   and performs aperture photometry
}{
   \sstdescription{
      This script reduces a ``nod jitter'' photometry observation with
      UFTI data.  It takes an UFTI observation comprising eight
      object frames and a dark frame to make a calibrated, untrimmed
      mosaic automatically.

      It performs a null debiassing, bad-pixel masking, dark
      subtraction, flat-field creation and division, feature
      detection and matching between object frames, and resampling.

      Photometry of the point source using a fixed 5-arcsecond aperture
      is calculated for each jitter frame and the mosaic.  The results
      appear in {\tt\$ORAC\_DATA\_OUT/aphot\_results.txt} in the form of a Starlink
      \xref{small text list}{sun190}{STLREF}.  The analysis of each star is appended to this file.

      This recipe works well for faint sources in moderately crowded fields.
   }
   \sstnotes{
      \sstitemlist{

         \sstitem
         The bad pixel mask applied is {\tt\$ORAC\_DATA\_CAL/bpm}.

         \sstitem
         The flat field is created by combining normalised object
         frames using the median at each pixel.

         \sstitem
         Registration is performed using common point sources in the
         overlap regions.  If the recipe cannot identify sufficient common
         objects, the script resorts to using the telescope offsets
         transformed to pixels.

         \sstitem
         The resampling applies non-integer shifts of origin using
         bilinear interpolation.  There is no rotation to align the
         Cartesian axes with the cardinal directions.

         \sstitem
         The recipe makes the mosaics by applying offsets in intensity
         to give the most consistent result amongst the overlapping regions.
         The mosaic is not trimmed to the dimensions of a single frame, thus
         the noise will be greater in the peripheral areas having received
         less exposure time.  The mosaic is not normalised by its exposure
         time (that being the exposure time of a single frame).

         \sstitem
         For each cycle of eight, the recipe creates a mosaic, which has
         its bad pixels filled and is then added into a master mosaic of
         improving signal to noise.  The exposure time is also summed and
         stored in the mosaic's EXP\_TIME header.  Likewise the end airmass
         header, AMEND, is updated to match that of the last-observed frame
         contributing to the mosaic.

         \sstitem
         Intermediate frames are deleted except for the flat-fielded ({\tt\_ff}
         suffix) frames.
      }
   }
   \sstdiytopic{
      Output Data
   }{
      \sstitemlist{

         \sstitem
         The integrated mosaic in {\tt g$<$date$>$\_$<$group\_number$>$\_mos}.

         \sstitem
         A mosaic for each cycle of eight in
         {\tt g$<$date$>$\_$<$group\_number$>$\_mos\_$<$cycle\_number$>$}, where {\tt$<$cycle\_number$>$}
         counts from 0.

         \sstitem
         The individual flat-fielded frames in {\tt f$<$date$>$\_$<$obs\_number$>$\_ff}.

         \sstitem
         The created flat fields in {\tt flat\_$<$group\_number$>$subgrp} for the
         first or only cycle, and {\tt flat\_$<$group\_number$>$subgrp\_$<$cycle\_number$>$}
         for subsequent cycles.
      }
   }
   \sstparameters{
      \sstsubsection{
         NUMBER = INTEGER
      }{
         The number of frames in the nod pattern.  It must be a multiple
         of four. {\tt[8]}
      }
   }
   \sstdiytopic{
      Related Recipes
   }{
      \htmlref{NOD4\_SELF\_FLAT\_NO\_MASK\_APHOT}{NOD4\_SELF\_FLAT\_NO\_MASK\_APHOT}.
   }
   \sstimplementationstatus{
      \sstitemlist{

         \sstitem
         The processing engines are from the Starlink packages: \xref{\CCDPACK}{sun139}{},
         \xref{\KAPPA}{sun95}{}, \xref{\PISA}{sun109}{}, and \xref{\PHOTOM}{sun45}{}.

         \sstitem
         Uses the Starlink NDF format.

         \sstitem
         History is recorded within the data files.

         \sstitem
         The title of the data is propagated through intermediate files
         to the mosaic.

         \sstitem
         Error propagation is not used.
      }
   }
}
\newpage
\sstroutine{
   QUADRANT\_JITTER
}{
   ``Quadrant Jitter'' standard reduction including object masking
}{
   \sstdescription{
      This script reduces a ``Quadrant Jitter'' photometry observation
      with UFTI data.  It takes an UFTI observation comprising one or
      more series of four object frames where the target is
      approximately centred in each quadrant; and a dark frame to make
      a calibrated, untrimmed mosaic automatically.

      It performs bad-pixel masking, null debiassing, dark subtraction,
      flat-field creation and division, feature detection and matching
      between object frames, and resampling.

      This recipe is suitable for faint objects or objects within
      a comparatively bright core embedded in faint extended emission,
      e.g. a quasar.  The object need not be isolated, as the recipe
      masks objects within the other quadrants, and hence does not
      introduce significant artifacts into the flat field.  For isolated
      objects use \htmlref{QUADRANT\_JITTER\_NO\_MASK}{QUADRANT\_JITTER\_NO\_MASK}; or where speed is critical,
      use QUADRANT\_JITTER\_BASIC instead.
   }
   \sstnotes{
      \sstitemlist{

         \sstitem
         The bad pixel mask applied is {\tt\$ORAC\_DATA\_CAL/bpm}.

         \sstitem
         The flat field is created iteratively.  First the quadrant
         containing the object is masked in each object frame.  Second an
         approximate flat field is created by combining the normalised
         and masked object frames using the clipped median at each pixel.
         This flat field is applied to the object frames.  Sources within
         the flat-fielded frames are detected, and masked in the
         dark-subtracted frames.  The second stage is repeated but applied
         to the masked frames to create the final flat field.

         \sstitem
         Registration is performed using common point sources in the
         overlap regions.  If the recipe cannot identify sufficient common
         objects, it matches the centroid of the central source.  If this
         fails, the script resorts to using the telescope offsets
         transformed to pixels.

         \sstitem
         The resampling applies non-integer shifts of origin using
         bilinear interpolation.  There is no rotation to align the
         Cartesian axes with the cardinal directions.

         \sstitem
         The recipe makes the mosaics by applying offsets in intensity
         to give the most consistent result amongst the overlapping regions.
         The mosaic is not trimmed to the dimensions of a single frame.  Thus
         the noise will be greater in the peripheral areas having received
         less exposure time.  The full signal will be in the central ninth
         containing the main object.  The mosaic is not normalised by its
         exposure time (that being the exposure time of a single frame).

         \sstitem
         For each cycle of four, the recipe creates a mosaic, which has
         its bad pixels filled and is then added into a master mosaic of
         improving signal to noise.  The exposure time is also summed and
         stored in the mosaic's EXP\_TIME header.  Likewise the end airmass
         header, AMEND, is updated to match that of the last-observed frame
         contributing to the mosaic.

         \sstitem
         Intermediate frames are deleted except for the flat-fielded ({\tt\_ff}
         suffix) frames.
      }
   }
   \sstdiytopic{
      Output Data
   }{
      \sstitemlist{

         \sstitem
         The integrated mosaic in {\tt g$<$date$>$\_$<$group\_number$>$\_mos}.

         \sstitem
         A mosaic for each cycle of four in
         {\tt g$<$date$>$\_$<$group\_number$>$\_mos\_$<$cycle\_number$>$}, where {\tt$<$cycle\_number$>$}
         counts from 0.

         \sstitem
         The individual flat-fielded frames in {\tt f$<$date$>$\_$<$obs\_number$>$\_ff}.

         \sstitem
         The created flat fields in {\tt flat\_$<$group\_number$>$subgrp} for the
         first or only cycle, and {\tt flat\_$<$group\_number$>$subgrp\_$<$cycle\_number$>$}
         for subsequent cycles.
      }
   }
   \sstdiytopic{
      Related Recipes
   }{
      \htmlref{QUADRANT\_JITTER\_BASIC}{QUADRANT\_JITTER\_BASIC},
      \htmlref{QUADRANT\_JITTER\_NO\_MASK}{QUADRANT\_JITTER\_NO\_MASK},\\
      QUADRANT\_JITTER\_TELE,\\
      \htmlref{EXTENDED\_5x5}{EXTENDED\_5x5},
      \htmlref{EXTENDED\_5x5\_BASIC}{EXTENDED\_5x5\_BASIC}.
   }
   \sstimplementationstatus{
      \sstitemlist{

         \sstitem
         The processing engines are from the Starlink packages: \xref{\CCDPACK}{sun139}{},
         \xref{\KAPPA}{sun95}{}, and  \xref{\PISA}{sun109}{}.

         \sstitem
         Uses the Starlink NDF format.

         \sstitem
         History is recorded within the data files.

         \sstitem
         The title of the data is propagated through intermediate files
         to the mosaic.

         \sstitem
         Error propagation is not used.
      }
   }
}
\newpage
\sstroutine{
   QUADRANT\_JITTER\_BASIC
}{
   ``Quadrant Jitter'' fastest version
}{
   \sstdescription{
      This script reduces a ``Quadrant Jitter'' photometry observation
      with UFTI data.  It takes an UFTI observation comprising one or
      more series of four object frames where the target is
      approximately centred in each quadrant; and a dark frame to make
      a calibrated, untrimmed mosaic automatically.

      It performs bad-pixel masking, null debiassing, dark subtraction,
      flat-field creation and division, and registration using telescope
      offsets.

      This recipe aims to keep pace with the pipeline's incoming
      data and many options which improve the final mosaic are omitted.
      This recipe is suitable for faint objects or objects within
      a comparatively bright core embedded in faint extended emission,
      e.g. a quasar.  If the object is not isolated, there will be
      artifacts introduced into the flat field.  This arise from the
      contribution of sources outside the quadrant containing the primary
      object.  This variant of \htmlref{QUADRANT\_JITTER}{QUADRANT\_JITTER} is best for isolated
      objects or where speed is critical.  Use \htmlref{QUADRANT\_JITTER}{QUADRANT\_JITTER} itself
      if object masking is required instead.
   }
   \sstnotes{
      \sstitemlist{

         \sstitem
         The bad pixel mask applied is {\tt\$ORAC\_DATA\_CAL/bpm}.

         \sstitem
         The flat field is created in two steps.  The quadrant
         containing the object is masked in each object frame.  Then the
         recipe combines the normalised and quadrant-masked object frames
         using the median at each pixel.

         \sstitem
         Registration is performed using the telescope offsets
         transformed to pixels.

         \sstitem
         There is no resampling, merely integer shifts of origin.

         \sstitem
         The recipe makes the mosaics by applying offsets in intensity
         to give the most consistent result amongst the overlapping regions.
         The mosaic is not trimmed to the dimensions of a single frame.  Thus
         the noise will be greater in the peripheral areas having received
         less exposure time.  The full signal will be in the central ninth
         containing the main object.  The mosaic is not normalised by its
         exposure time (that being the exposure time of a single frame).

         \sstitem
         For each cycle of four, the recipe creates a mosaic, which is
         added into a master mosaic of improving signal to noise.  The
         exposure time is also summed and stored in the mosaic's EXP\_TIME
         header.  Likewise the end airmass header, AMEND, is updated to
         match that of the last-observed frame contributing to the mosaic.

         \sstitem
         Intermediate frames are deleted except for the flat-fielded ({\tt\_ff}
         suffix) frames.
      }
   }
   \sstdiytopic{
      Output Data
   }{
      \sstitemlist{

         \sstitem
         The integrated mosaic in {\tt g$<$date$>$\_$<$group\_number$>$\_mos}.

         \sstitem
         A mosaic for each cycle of four in
         {\tt g$<$date$>$\_$<$group\_number$>$\_mos\_$<$cycle\_number$>$}, where {\tt$<$cycle\_number$>$}
         counts from 0.

         \sstitem
         The individual flat-fielded frames in {\tt f$<$date$>$\_$<$obs\_number$>$\_ff}.

         \sstitem
         The created flat fields in {\tt flat\_$<$group\_number$>$subgrp} for the
         first or only cycle, and {\tt flat\_$<$group\_number$>$subgrp\_$<$cycle\_number$>$}
         for subsequent cycles.
      }
   }
   \sstdiytopic{
      Related Recipes
   }{
      \htmlref{QUADRANT\_JITTER}{QUADRANT\_JITTER},
      \htmlref{QUADRANT\_JITTER\_NO\_MASK}{QUADRANT\_JITTER\_NO\_MASK},
      \htmlref{EXTENDED\_5x5}{EXTENDED\_5x5}.
   }
   \sstimplementationstatus{
      \sstitemlist{

         \sstitem
         The processing engines are from the Starlink packages: \xref{\CCDPACK}{sun139}{},
         \xref{\KAPPA}{sun95}{}, and  \xref{\PISA}{sun109}{}.

         \sstitem
         Uses the Starlink NDF format.

         \sstitem
         History is recorded within the data files.

         \sstitem
         The title of the data is propagated through intermediate files
         to the mosaic.

         \sstitem
         Error propagation is not used.
      }
   }
}
\newpage
\sstroutine{
   QUADRANT\_JITTER\_NO\_MASK
}{
   ``Quadrant Jitter'' standard reduction without object masking
}{
   \sstdescription{
      This script reduces a ``Quadrant Jitter'' photometry observation
      with UFTI data.  It takes an UFTI observation comprising one or
      more series of four object frames where the target is
      approximately centred in each quadrant; and a dark frame to make
      a calibrated, untrimmed mosaic automatically.

      It performs bad-pixel masking, null debiassing, dark subtraction,
      flat-field creation and division, feature detection and matching
      between object frames, and resampling.

      This recipe is suitable for faint objects or objects within
      a comparatively bright core embedded in faint extended emission,
      e.g. a quasar.  If there are other objects of comparable
      brightness to the principal target in other quadrants, they will
      introduce artifacts into the flat field.  Use \htmlref{QUADRANT\_JITTER}{QUADRANT\_JITTER} for
      those.
   }
   \sstnotes{
      \sstitemlist{

         \sstitem
         The bad pixel mask applied is {\tt\$ORAC\_DATA\_CAL/bpm}.

         \sstitem
         The flat field is created in two steps.  The quadrant
         containing the object is masked in each object frame.  Then the
         recipe combines the normalised and quadrant-masked object frames
         using the median at each pixel.

         \sstitem
         Registration is performed using common point sources in the
         overlap regions.  If the recipe cannot identify sufficient common
         objects, it matches the centroid of the central source.  If this
         fails, the script resorts to using the telescope offsets
         transformed to pixels.

         \sstitem
         The resampling applies non-integer shifts of origin using
         bilinear interpolation.  There is no rotation to align the
         Cartesian axes with the cardinal directions.

         \sstitem
         The recipe makes the mosaics by applying offsets in intensity
         to give the most consistent result amongst the overlapping regions.
         The mosaic is not trimmed to the dimensions of a single frame.  Thus
         the noise will be greater in the peripheral areas having received
         less exposure time.  The full signal will be in the central ninth
         containing the main object.  The mosaic is not normalised by its
         exposure time (that being the exposure time of a single frame).

         \sstitem
         For each cycle of four, the recipe creates a mosaic, which has
         its bad pixels filled and is then added into a master mosaic of
         improving signal to noise.  The exposure time is also summed and
         stored in the mosaic's EXP\_TIME header.  Likewise the end airmass
         header, AMEND, is updated to match that of the last-observed frame
         contributing to the mosaic.

         \sstitem
         Intermediate frames are deleted except for the flat-fielded ({\tt\_ff}
         suffix) frames.
      }
   }
   \sstdiytopic{
      Output Data
   }{
      \sstitemlist{

         \sstitem
         The integrated mosaic in {\tt g$<$date$>$\_$<$group\_number$>$\_mos}.

         \sstitem
         A mosaic for each cycle of four in
         {\tt g$<$date$>$\_$<$group\_number$>$\_mos\_$<$cycle\_number$>$}, where {\tt$<$cycle\_number$>$}
         counts from 0.

         \sstitem
         The individual flat-fielded frames in {\tt f$<$date$>$\_$<$obs\_number$>$\_ff}.

         \sstitem
         The created flat fields in {\tt flat\_$<$group\_number$>$subgrp} for the
         first or only cycle, and {\tt flat\_$<$group\_number$>$subgrp\_$<$cycle\_number$>$}
         for subsequent cycles.
      }
   }
   \sstdiytopic{
      Related Recipes
   }{
      \htmlref{QUADRANT\_JITTER}{QUADRANT\_JITTER},
      \htmlref{QUADRANT\_JITTER\_BASIC}{QUADRANT\_JITTER\_BASIC},
      QUADRANT\_JITTER\_TELE,\\
      \htmlref{EXTENDED\_5x5}{EXTENDED\_5x5}, 
      \htmlref{EXTENDED\_5x5\_BASIC}{EXTENDED\_5x5\_BASIC}.
   }
   \sstimplementationstatus{
      \sstitemlist{

         \sstitem
         The processing engines are from the Starlink packages: \xref{\CCDPACK}{sun139}{}
         and \xref{\KAPPA}{sun95}{}.

         \sstitem
         Uses the Starlink NDF format.

         \sstitem
         History is recorded within the data files.

         \sstitem
         The title of the data is propagated through intermediate files
         to the mosaic.

         \sstitem
         Error propagation is not used.
      }
   }
}
\newpage
\sstroutine{
   REDUCE\_DARK
}{
   File an observation as the current dark
}{
   \sstdescription{
      This recipe reduces a dark frame observation with UFTI
      data.  It files the dark in file {\tt\$ORAC\_DATA\_OUT/index.dark}.
      Reduction is only bad-pixel masking.
   }
   \sstnotes{
      \sstitemlist{

         \sstitem
         The bad pixel mask applied is {\tt\$ORAC\_DATA\_CAL/bpm}.

         \sstitem
         Intermediate frames are deleted.

         \sstitem
         Sub-arrays are supported.
      }
   }
   \sstdiytopic{
      Output Data
   }{
      \sstitemlist{

         \sstitem
         The dark is called {\tt dark\_$<$frame\_number$>$}.

         \sstitem
         The dark is filed in {\tt f$<$\$ORAC\_DATA\_OUT/index.dark$>$}.
      }
   }
   \sstdiytopic{
      Timing
   }{
      Of the order of a few seconds on an Ultra 5.
   }
   \sstimplementationstatus{
      \sstitemlist{

         \sstitem
         The processing engines are from the Starlink package \xref{\KAPPA}{sun95}{}.

         \sstitem
         Uses the Starlink NDF format.

         \sstitem
         History is recorded within the data files.

         \sstitem
         The title of the data is propagated through the intermediate file
         to the dark.

         \sstitem
         Error propagation is not used.
      }
   }
}
\newpage
\sstroutine{
   SKY\_AND\_JITTER5
}{
   Reduces a 5-point ``combined jitter'' photometry observation
}{
   \sstdescription{
      This script reduces a ``combined jitter'' photometry observation
      with UFTI data.  It takes an UFTI observation comprising a one or
      more sets of frames, each set containing a sky frame, followed by
      five object frames; and a pre-determined flat-field frame to make
      a calibrated, trimmed mosaic automatically.

      This recipe performs bad-pixel masking, null debiassing, sky
      subtraction, flat-field division, feature detection and matching
      between object frames, and resampling.  See the notes for details.

      This recipe is suitable for moderately faint point sources.
   }
   \sstnotes{
      \sstitemlist{

         \sstitem
         You may use \htmlref{SKY\_FLAT}{SKY\_FLAT} or \htmlref{SKY\_FLAT\_MASKED}{SKY\_FLAT\_MASKED} to make the flat field.

         \sstitem
         The bad pixel mask applied is {\tt\$ORAC\_DATA\_CAL/bpm}.

         \sstitem
         The most-recent sky frame is used for the sky subtraction.

         \sstitem
         Where automatic registration is not possible, the recipe matches
         the centroid of central source, and should that fail, it resorts
         to using the telescope offsets transformed to pixels.

         \sstitem
         The resampling applies non-integer shifts of origin using
         bilinear interpolation.  There is no rotation to align the
         Cartesian axes with the cardinal directions.

         \sstitem
         The recipe makes the mosaic by applying offsets in intensity to
         give the most consistent result amongst the overlapping regions.
         The mosaic is trimmed to the dimensions of an input frame.  The
         mosaic is not normalised by its exposure time (that being the
         exposure time of a single frame).

         \sstitem
         For each cycle of five, the recipe creates a mosaic, which
         is then added into a master mosaic of improving signal to noise.
         The exposure time is also summed and stored in the mosaic's
         EXP\_TIME header.  Likewise the end airmass header, AMEND, is
         updated to match that of the last-observed frame contributing to
         the mosaic.

         \sstitem
         Intermediate frames are deleted except for the flat-fielded ({\tt\_ff}
         suffix) frames.

         \sstitem
         Sub-arrays are supported.
      }
   }
   \sstdiytopic{
      Output Data
   }{
      \sstitemlist{

         \sstitem
         The resultant mosaic in {\tt g$<$date$>$\_$<$group\_number$>$\_mos}.

         \sstitem
         The individual flat-fielded frames in {\tt f$<$date$>$\_$<$obs\_number$>$\_ff}.
      }
   }
   \sstparameters{
      \sstsubsection{
         NUMBER = INTEGER
      }{
         The number of frames in the jitter pattern. {\tt[5]}
      }
   }
   \sstdiytopic{
      Related Recipes
   }{
      \htmlref{SKY\_AND\_JITTER5\_APHOT}{SKY\_AND\_JITTER5\_APHOT},
      \htmlref{SKY\_FLAT}{SKY\_FLAT},
      \htmlref{SKY\_FLAT\_MASKED}{SKY\_FLAT\_MASKED}.
   }
   \sstimplementationstatus{
      \sstitemlist{

         \sstitem
         The processing engines are from the Starlink packages: \xref{\CCDPACK}{sun139}{}
         and \xref{\KAPPA}{sun95}{}.

         \sstitem
         Uses the Starlink NDF format.

         \sstitem
         History is recorded within the data files.

         \sstitem
         The title of the data is propagated through intermediate files
         to the mosaic.

         \sstitem
         Error propagation is not used.
      }
   }
}
\newpage
\sstroutine{
   SKY\_AND\_JITTER5\_APHOT
}{
   Reduces a 5-point ``combined jitter'' photometry observation, and
   performs aperture photometry
}{
   \sstdescription{
      This script reduces a ``combined jitter'' photometry observation
      with UFTI data.  It takes an UFTI observation comprising a one or
      more sets of frames, each set containing a sky frame, followed by
      five object frames; and a pre-determined flat-field frame to make
      a calibrated, trimmed mosaic automatically.

      This recipe performs bad-pixel masking, null debiassing, sky
      subtraction, flat-field division, feature detection and matching
      between object frames, and resampling.  See the notes for details.

      Photometry of the point source using a fixed 6-arcsecond aperture
      is calculated for each jitter frame and the mosaic.  The results
      appear in {\tt\$ORAC\_DATA\_OUT/aphot\_results.txt} in the form of
      a Starlink
      \xref{small text list}{sun190}{STLREF}.  The analysis of each star is appended to this file.

      This recipe is suitable for moderately faint point sources.
   }
   \sstnotes{
      \sstitemlist{

         \sstitem
         You may use \htmlref{SKY\_FLAT}{SKY\_FLAT} or \htmlref{SKY\_FLAT\_MASKED}{SKY\_FLAT\_MASKED} to make the flat field.

         \sstitem
         The bad pixel mask applied is {\tt\$ORAC\_DATA\_CAL/bpm}.

         \sstitem
         The most-recent sky frame is used for the sky subtraction.

         \sstitem
         Where automatic registration is not possible, the recipe matches
         the centroid of central source, and should that fail, it resorts
         to using the telescope offsets transformed to pixels.

         \sstitem
         The resampling applies non-integer shifts of origin using
         bilinear interpolation.  There is no rotation to align the
         Cartesian axes with the cardinal directions.

         \sstitem
         The recipe makes the mosaic by applying offsets in intensity to
         give the most consistent result amongst the overlapping regions.
         The mosaic is trimmed to the dimensions of an input frame.  The
         mosaic is not normalised by its exposure time (that being the
         exposure time of a single frame).

         \sstitem
         For each cycle of five, the recipe creates a mosaic, which
         is then added into a master mosaic of improving signal to noise.
         The exposure time is also summed and stored in the mosaic's
         EXP\_TIME header.  Likewise the end airmass header, AMEND, is
         updated to match that of the last-observed frame contributing to
         the mosaic.

         \sstitem
         Intermediate frames are deleted except for the flat-fielded ({\tt\_ff}
         suffix) frames.

         \sstitem
         Sub-arrays are supported.
      }
   }
   \sstdiytopic{
      Output Data
   }{
      \sstitemlist{

         \sstitem
         The resultant mosaic in {\tt g$<$date$>$\_$<$group\_number$>$\_mos}.

         \sstitem
         The individual flat-fielded frames in {\tt f$<$date$>$\_$<$obs\_number$>$\_ff}.
      }
   }
   \sstparameters{
      \sstsubsection{
         NUMBER = INTEGER
      }{
         The number of frames in the jitter pattern. {\tt[5]}
      }
   }
   \sstdiytopic{
      Related Recipes
   }{
      \htmlref{JITTER5\_SELF\_FLAT\_APHOT}{JITTER5\_SELF\_FLAT\_APHOT},
      \htmlref{SKY\_AND\_JITTER5}{SKY\_AND\_JITTER5},
      \htmlref{SKY\_FLAT}{SKY\_FLAT},\\
      \htmlref{SKY\_FLAT\_MASKED}{SKY\_FLAT\_MASKED}.
   }
   \sstimplementationstatus{
      \sstitemlist{

         \sstitem
         The processing engines are from the Starlink packages: \xref{\CCDPACK}{sun139}{},
         \xref{\KAPPA}{sun95}{}, and \xref{\PHOTOM}{sun45}{}.

         \sstitem
         Uses the Starlink NDF format.

         \sstitem
         History is recorded within the data files.

         \sstitem
         The title of the data is propagated through intermediate files
         to the mosaic.

         \sstitem
         Error propagation is not used.
      }
   }
}
\newpage
\sstroutine{
   SKY\_FLAT
}{
   Creates and files a flat field derived from 5 jittered frames
}{
   \sstdescription{
      This recipe make a sky flat from a series of object frames (these may
      include blank sky) combined using one of a selection of statistics.
      It expects one dark frame followed by five object frames.

      It performs a null debiassing, bad-pixel masking, and dark subtraction
      before combining normalised frames pixel by pixel using the median.
      Details of the flat is filed in the index of flats for future
      selection and use of the flat.

      For best results the field observed should contain few stars and no
      bright ones.  In contaminated sky regions, recipe 
      \htmlref{SKY\_FLAT\_MASKED}{SKY\_FLAT\_MASKED} will
      greatly reduce artifacts appearing in the resultant flat.
   }
   \sstnotes{
      \sstitemlist{

         \sstitem
         The bad pixel mask applied is \$ORAC\_DATA\_CAL/bpm.

         \sstitem
         Intermediate frames are deleted.

         \sstitem
         Sub-arrays are supported.
      }
   }
   \sstdiytopic{
      Output Data
   }{
      \sstitemlist{

         \sstitem
         The flat is called {\tt flat\_$<$group\_number$>$}.

         \sstitem
         The flat is filed in {\tt $<$\$ORAC\_DATA\_OUT/index.flat$>$}.
      }
   }
   \sstdiytopic{
      Related Recipes
   }{
      \htmlref{SKY\_FLAT\_MASKED}{SKY\_FLAT\_MASKED}.
   }
   \sstimplementationstatus{
      \sstitemlist{

         \sstitem
         The processing engines are from the Starlink packages
         \xref{\CCDPACK}{sun139}{} and \xref{\KAPPA}{sun95}{}.

         \sstitem
         Uses the Starlink NDF format.

         \sstitem
         History is recorded within the data files.

         \sstitem
         The title of the data is propagated through the intermediate file
         to the flat.

         \sstitem
         Error propagation is not used.
      }
   }
}
\newpage
\sstroutine{
   SKY\_FLAT\_MASKED
}{
   Creates and files a flat field derived from 5 jittered frames using
   object masking to reduce artifacts
}{
   \sstdescription{
      This recipe make a sky flat from a series of object frames (these may
      include blank sky) combined using one of a selection of statistics.
      It expects one dark frame followed by five object frames.

      It performs a null debiassing, bad-pixel masking, and dark subtraction
      before combining the object frames frames pixel by pixel to make the flat.
      See the notes for further details.  The parameters of the flat is filed
      in the index of flats for future selection and use of the flat.

      For best results the field observed should contain few stars and no
      bright ones.  In contaminated sky regions, recipe \htmlref{SKY\_FLAT\_MASKED}{SKY\_FLAT\_MASKED} will
      greatly reduce artifacts appearing in the resultant flat.
   }
   \sstnotes{
      \sstitemlist{

         \sstitem
         The bad pixel mask applied is {\tt\$ORAC\_DATA\_CAL/bpm}.

         \sstitem
         The flat field is created iteratively.  First an approximate
         flat-field is created by combining normalised object frames using
         the median at each pixel.  This flat field is applied to the object
         frames.  Sources within the flat-fielded frames are detected, and
         masked in the dark-subtracted frames.  The first stage is repeated
         but applied to the masked frames to create the final flat field.

         \sstitem
         Intermediate frames are deleted.

         \sstitem
         Sub-arrays are supported.
      }
   }
   \sstdiytopic{
      Output Data
   }{
      \sstitemlist{

         \sstitem
         The flat is called {\tt flat\_$<$group\_number$>$}.

         \sstitem
         The flat is filed in {\tt $<$\$ORAC\_DATA\_OUT/index.flat$>$}.
      }
   }
   \sstdiytopic{
      Related Recipes
   }{
      \htmlref{SKY\_FLAT}{SKY\_FLAT}.
   }
   \sstimplementationstatus{
      \sstitemlist{

         \sstitem
         The processing engines are from the Starlink packages \xref{\CCDPACK}{sun139}{},
         \xref{\KAPPA}{sun95}{}, and \xref{\PISA}{sun109}{}.

         \sstitem
         Uses the Starlink NDF format.

         \sstitem
         History is recorded within the data files.

         \sstitem
         The title of the data is propagated through the intermediate file
         to the flat.

         \sstitem
         Error propagation is not used.
      }
   }
}

% ? End of main text
\end{document}
